\documentclass[aps,preprint]{revtex4}%
\usepackage{amssymb}
\usepackage{amsfonts}
\usepackage{amsmath}
\usepackage{graphicx}%
\setcounter{MaxMatrixCols}{30}
%TCIDATA{OutputFilter=latex2.dll}
%TCIDATA{Version=5.50.0.2960}
%TCIDATA{CSTFile=revtex4.cst}
%TCIDATA{Created=Sunday, March 14, 2010 09:37:54}
%TCIDATA{LastRevised=Monday, June 15, 2015 10:20:46}
%TCIDATA{<META NAME="GraphicsSave" CONTENT="32">}
%TCIDATA{<META NAME="SaveForMode" CONTENT="1">}
%TCIDATA{BibliographyScheme=Manual}
%TCIDATA{<META NAME="DocumentShell" CONTENT="Articles\SW\REVTeX 4">}
%BeginMSIPreambleData
\providecommand{\U}[1]{\protect\rule{.1in}{.1in}}
%EndMSIPreambleData
\newtheorem{theorem}{Theorem}
\newtheorem{acknowledgement}[theorem]{Acknowledgement}
\newtheorem{algorithm}[theorem]{Algorithm}
\newtheorem{axiom}[theorem]{Axiom}
\newtheorem{claim}[theorem]{Claim}
\newtheorem{conclusion}[theorem]{Conclusion}
\newtheorem{condition}[theorem]{Condition}
\newtheorem{conjecture}[theorem]{Conjecture}
\newtheorem{corollary}[theorem]{Corollary}
\newtheorem{criterion}[theorem]{Criterion}
\newtheorem{definition}[theorem]{Definition}
\newtheorem{example}[theorem]{Example}
\newtheorem{exercise}[theorem]{Exercise}
\newtheorem{lemma}[theorem]{Lemma}
\newtheorem{notation}[theorem]{Notation}
\newtheorem{problem}[theorem]{Problem}
\newtheorem{proposition}[theorem]{Proposition}
\newtheorem{remark}[theorem]{Remark}
\newtheorem{solution}[theorem]{Solution}
\newtheorem{summary}[theorem]{Summary}
\newenvironment{proof}[1][Proof]{\noindent\textbf{#1.} }{\ \rule{0.5em}{0.5em}}
\begin{document}
\preprint{ }
\title[ ]{REV\TeX4}
\author{Brian Weiner}
\affiliation{}
\author{}
\affiliation{Penn State}
\author{}
\startpage{1}
\endpage{ }
\maketitle
\tableofcontents


\section{Procedure}

\begin{itemize}
\item Start with a geminal on an arbitary basis

\item Symmetrically orthonormalize the basis (Not unique if eigenvalues of
metric matrix is degenerate)

\item Form the geminal FORDO

\item Diagonalize this operator

\item Create the antisymmetric operator (The geminal is invariant to
$USp\left(  2s,\mathbb{C}\right)  $ and this group will generate all possible
conb representations)

\item Find canonical form (Invariant to $SU(1,1)^{s}$ for non degenerate
geminals and $SU(1,1)^{s-p}\times SU(p_{1},p_{1})\times\cdots\times
SU(p_{\kappa},p_{\kappa})$ for degenerate ones)
\end{itemize}

\section{Geminals}

Every geminal $\left\vert g\right\rangle \in\mathcal{H}^{2}$ defines a
symplectic form on $\mathcal{H}^{1}$ given by
\begin{equation}
g\left(  \varphi_{1},\varphi_{2}\right)  =\left\langle \left.  g\right\vert
\varphi_{1}\wedge\varphi_{2}\right\rangle =-g\left(  \varphi_{2},\varphi
_{1}\right)  ;\left\vert \varphi_{1}\right\rangle ,\left\vert \varphi
_{2}\right\rangle \in\mathcal{H}^{1}%
\end{equation}
that can be expressed in canonical form as
\begin{equation}
\left\vert g\right\rangle =\sum_{1\leq j\leq s}c_{j}\left\vert \psi_{j}%
\psi_{-j}\right\rangle =\sum_{1\leq j\leq s}n_{j}^{\frac{1}{2}}\left\vert
\psi_{j}\psi_{-j}\right\rangle =\sum_{1\leq j\leq s}\left\vert c_{j}^{\frac
{1}{2}}\psi_{j}c_{j}^{\frac{1}{2}}\psi_{-j}\right\rangle =\sum_{1\leq j\leq
s}\left\vert \eta_{j}\eta_{-j}\right\rangle ,
\end{equation}
where $\left\{  \left.  \left\vert \psi_{j}\right\rangle \right\vert -s\leq
j\leq s\right\}  $ is a complete orthonormal basis for $\mathcal{H}^{1}$ and
$c_{j}\geq0,1\leq j\leq s.$(The pairs of GSO's $\left\vert \psi_{j}%
\right\rangle ,\left\vert \psi_{-j}\right\rangle $ are not unique as
$\left\vert \psi_{j}\psi_{-j}\right\rangle $ are invariant under $SU(2)$).

The resolution of the identity on $\mathcal{H}^{2}$ can be expressed as
\begin{align}
I_{2}  &  =\sum_{\substack{1\leq j<k\leq r\\1\leq l<m\leq r}}\left\vert
\varphi_{j}\varphi_{k}\right\rangle \left[  \mathbf{\Lambda}^{\left(
2\right)  }\left(  \Delta_{\varphi}\right)  \right]  _{jklm}^{-1}\left\langle
\varphi_{l}\varphi_{m}\right\vert \nonumber\\
&  =\sum_{\substack{1\leq j<k\leq r\\1\leq l<m\leq r\\1\leq p<q\leq
r}}\left\vert \varphi_{j}\varphi_{k}\right\rangle \left[  \mathbf{\Lambda
}^{\left(  2\right)  }\left(  \Delta_{\varphi}\right)  \right]  _{jklm}%
^{-\frac{1}{2}}\left[  \mathbf{\Lambda}^{\left(  2\right)  }\left(
\Delta_{\varphi}\right)  \right]  _{lmpq}^{-\frac{1}{2}}\left\langle
\varphi_{p}\varphi_{q}\right\vert ,\nonumber\\
&  =\sum_{1\leq j<k\leq r}\left\vert \tilde{\varphi}_{j}\tilde{\varphi}%
_{k}\right\rangle \left\langle \tilde{\varphi}_{j}\tilde{\varphi}%
_{k}\right\vert
\end{align}
producing
\begin{equation}
\sum_{1\leq j<k\leq r}\left\vert \varphi_{j}\varphi_{k}\right\rangle \left[
\mathbf{\Lambda}^{\left(  2\right)  }\left(  \Delta_{\varphi}\right)  \right]
_{jklm}^{-1}\left\langle \left.  \varphi_{l}\varphi_{m}\right\vert
g\right\rangle =\sum_{1\leq j<k\leq s}\left\vert \varphi_{j}\varphi
_{k}\right\rangle \left(  \left[  \mathbf{\Lambda}^{\left(  2\right)  }\left(
\Delta_{\varphi}\right)  \right]  ^{-1}\mathbf{g}_{2\varphi}\right)  _{jk}%
\end{equation}
and
\begin{equation}
\sum_{1\leq j<k\leq r}\left\vert \tilde{\varphi}_{j}\tilde{\varphi}%
_{k}\right\rangle \left\langle \left.  \tilde{\varphi}_{j}\tilde{\varphi}%
_{k}\right\vert g\right\rangle =\sum_{1\leq j<k\leq r}\left\vert
\tilde{\varphi}_{j}\tilde{\varphi}_{k}\right\rangle \mathbf{g}_{2\tilde
{\varphi}jk}.
\end{equation}
This gives that
\begin{equation}
\mathbf{g}_{2\tilde{\varphi}}=\left[  \mathbf{\Lambda}^{\left(  2\right)
}\left(  \Delta_{\varphi}\right)  \right]  ^{-\frac{1}{2}}\mathbf{g}%
_{2\varphi}.
\end{equation}
Thus the expansion of a geminal wrt a general basis $\left\{  \left.
\left\vert \varphi_{j}\right\rangle \right\vert 1\leq j\leq2s=r\right\}  $ is
\begin{align}
\left\vert g\right\rangle  &  =\frac{1}{2}\sum_{1\leq j,k\leq s}\left\vert
\varphi_{j}\varphi_{k}\right\rangle \mathcal{R}_{2\varphi}\left(  g\right)
_{jk}\nonumber\\
&  =\frac{1}{2}\sum_{1\leq j,k\leq s}\left\vert \varphi_{j}\varphi
_{k}\right\rangle \left(  \left[  \mathbf{\Lambda}^{\left(  2\right)  }\left(
\Delta_{\varphi}\right)  \right]  ^{-1}\mathbf{g}_{2\varphi}\right)  _{jk},
\end{align}
$\lceil$%
\begin{align}
\mathcal{R}_{2\varphi}\left(  g\right)  _{jk}  &  =-\mathcal{R}_{2\varphi
}\left(  g\right)  _{kj}\nonumber\\
\left(  \left[  \mathbf{\Lambda}^{\left(  2\right)  }\left(  \Delta_{\varphi
}\right)  \right]  ^{-1}\mathbf{g}_{2\varphi}\right)  _{jk}  &  =-\left(
\left[  \mathbf{\Lambda}^{\left(  2\right)  }\left(  \Delta_{\varphi}\right)
\right]  ^{-1}\mathbf{g}_{2\varphi}\right)  _{kj}\rfloor
\end{align}
where
\begin{equation}
\left\vert g\right\rangle =\left(  \left\vert \varphi_{1}\wedge\varphi
_{2}\right\rangle ,\cdots,\left\vert \varphi_{r-1}\wedge\varphi_{r}%
\right\rangle \right)  \mathcal{R}_{2\varphi}\left(  g\right)
\end{equation}
and
\begin{align}
\mathbf{g}_{2\varphi}  &  =\left\langle \left.  \left(
\begin{array}
[c]{c}%
\left\langle \varphi_{1}\wedge\varphi_{2}\right\vert \\
\vdots\\
\left\langle \varphi_{r-1}\wedge\varphi_{r}\right\vert
\end{array}
\right)  \right\vert g\right\rangle \\
&  =\left\langle \left.  \left(
\begin{array}
[c]{c}%
\left\langle \varphi_{1}\wedge\varphi_{2}\right\vert \\
\vdots\\
\left\langle \varphi_{r-1}\wedge\varphi_{r}\right\vert
\end{array}
\right)  \right\vert \left(  \left\vert \varphi_{1}\wedge\varphi
_{2}\right\rangle ,\cdots,\left\vert \varphi_{r-1}\wedge\varphi_{r}%
\right\rangle \right)  \right\rangle \mathcal{R}_{2\varphi}\left(  g\right) \\
&  =\mathbf{\Lambda}^{\left(  2\right)  }\left(  \Delta_{\varphi}\right)
\mathcal{R}_{2\varphi}\left(  g\right)  .
\end{align}
In the orthonormalized basis
\begin{equation}
\mathcal{R}_{2\tilde{\varphi}}\left(  g\right)  =\mathbf{g}_{2\tilde{\varphi}}%
\end{equation}
and in the general basis
\begin{equation}
\mathcal{R}_{2\varphi}\left(  g\right)  =\left[  \mathbf{\Lambda}^{\left(
2\right)  }\left(  \Delta_{\varphi}\right)  \right]  ^{-1}\mathbf{g}%
_{2\varphi}%
\end{equation}
thus
\begin{equation}
\mathcal{R}_{2\tilde{\varphi}}\left(  g\right)  =\left[  \mathbf{\Lambda
}^{\left(  2\right)  }\left(  \Delta_{\varphi}\right)  \right]  ^{\frac{1}{2}%
}\mathcal{R}_{2\varphi}\left(  g\right)
\end{equation}
as%
\begin{equation}
\mathcal{R}_{2\tilde{\varphi}}\left(  g\right)  =\mathbf{g}_{2\tilde{\varphi}%
}=\left[  \mathbf{\Lambda}^{\left(  2\right)  }\left(  \Delta_{\varphi
}\right)  \right]  ^{-\frac{1}{2}}\mathbf{g}_{2\varphi}.
\end{equation}


\subsubsection{Geminal FORDO}

Each set $\left\{  e^{i\theta}\left\vert g\right\rangle \right\}  $ of
normalized geminals can be associated with a projection operator $\left\vert
g\right\rangle \left\langle g\right\vert ,$ that has the representation%
\begin{align}
&  \left\vert g\right\rangle \left\langle g\right\vert =\\
&  \sum_{\substack{1\leq j_{1}<k_{1}\leq r\\1\leq l_{1}<m_{1}\leq r\\1\leq
j_{2}<k_{2}\leq r\\1\leq l_{2}<m_{2}\leq r}}\left\vert \varphi_{j_{1}}%
\varphi_{k_{1}}\right\rangle \left[  \mathbf{\Lambda}^{\left(  2\right)
}\left(  \Delta_{\varphi}\right)  \right]  _{j_{1}k_{1}l_{1}m_{1}}%
^{-1}\left\langle \varphi_{l_{1}}\varphi_{m_{1}}\right\vert \left\vert
g\right\rangle \left\langle g\right\vert \left\vert \varphi_{j_{2}}%
\varphi_{k_{2}}\right\rangle \left[  \mathbf{\Lambda}^{\left(  2\right)
}\left(  \Delta_{\varphi}\right)  \right]  _{j_{2}k_{2}l_{2}m_{2}}%
^{-1}\left\langle \varphi_{l_{2}}\varphi_{m_{2}}\right\vert \\
&  \sum_{\substack{1\leq j_{1}<k_{1}\leq r\\1\leq l_{1}<m_{1}\leq r\\1\leq
j_{2}<k_{2}\leq r\\1\leq l_{2}<m_{2}\leq r}}\left\vert \varphi_{j_{1}}%
\varphi_{k_{1}}\right\rangle \left[  \mathbf{\Lambda}^{\left(  2\right)
}\left(  \Delta_{\varphi}\right)  \right]  _{j_{1}k_{1}l_{1}m_{1}}%
^{-1}g_{2\varphi l_{1}m_{1}}\bar{g}_{2\varphi j_{2}k_{2}}\left[
\mathbf{\Lambda}^{\left(  2\right)  }\left(  \Delta_{\varphi}\right)  \right]
_{j_{2}k_{2}l_{2}m_{2}}^{-1}\left\langle \varphi_{l_{2}}\varphi_{m_{2}%
}\right\vert \\
&  \sum_{\substack{1\leq j_{1}<k_{1}\leq r\\1\leq l_{2}<m_{2}\leq
r}}\left\vert \varphi_{j_{1}}\varphi_{k_{1}}\right\rangle \left(  \left[
\mathbf{\Lambda}^{\left(  2\right)  }\left(  \Delta_{\varphi}\right)  \right]
^{-1}\mathbf{g}_{2\varphi}\otimes\mathbf{g}_{2\varphi}^{\dagger}\left[
\mathbf{\Lambda}^{\left(  2\right)  }\left(  \Delta_{\varphi}\right)  \right]
^{-1}\right)  _{j_{1}k_{1}l_{2}m_{2}}\left\langle \varphi_{l_{2}}%
\varphi_{m_{2}}\right\vert \\
&  \sum_{\substack{1\leq j_{1}<k_{1}\leq r\\1\leq l_{2}<m_{2}\leq
r}}\left\vert \varphi_{j_{1}}\varphi_{k_{1}}\right\rangle \left(
\mathcal{R}_{2\varphi}\left(  \left\vert g\right\rangle \left\langle
g\right\vert \right)  \otimes\mathcal{R}_{2\varphi}\left(  \left\vert
g\right\rangle \left\langle g\right\vert \right)  ^{\dagger}\right)
_{j_{1}k_{1}l_{2}m_{2}}\left\langle \varphi_{l_{2}}\varphi_{m_{2}}\right\vert
\end{align}
the contraction of this projection operator is the FORDO of the geminal given
by
\begin{equation}
2L_{2}^{1}\left(  \left\vert g\right\rangle \left\langle g\right\vert \right)
=D^{1}\left(  g\right)  ,
\end{equation}
which has the representation%
\begin{align}
&  \frac{1}{4}\sum_{\substack{1\leq j_{1},k_{1}\leq r\\1\leq l_{2},m_{2}\leq
r}}\left\vert \varphi_{j_{1}}\varphi_{k_{1}}\right\rangle \left(  \left[
\mathbf{\Lambda}^{\left(  2\right)  }\left(  \Delta_{\varphi}\right)  \right]
^{-\frac{1}{2}}\left[  \mathbf{\Lambda}^{\left(  2\right)  }\left(
\Delta_{\varphi}\right)  \right]  ^{-\frac{1}{2}}\mathbf{g}_{2\varphi}%
\otimes\mathbf{g}_{2\varphi}^{\dagger}\left[  \mathbf{\Lambda}^{\left(
2\right)  }\left(  \Delta_{\varphi}\right)  \right]  \right.  ^{-\frac{1}{2}%
}\\
&  \left.  \left[  \mathbf{\Lambda}^{\left(  2\right)  }\left(  \Delta
_{\varphi}\right)  \right]  ^{-\frac{1}{2}}\right)  _{j_{1}k_{1}l_{2}m_{2}%
}\left\langle \varphi_{l_{2}}\varphi_{m_{2}}\right\vert \\
&  =\frac{1}{4}\sum_{\substack{1\leq j_{1},k_{1}\leq r\\1\leq l_{2},m_{2}\leq
r}}\left\vert \left(  \Delta_{\varphi}^{-\frac{1}{2}}\varphi\right)  _{j_{1}%
}\left(  \Delta_{\varphi}^{-\frac{1}{2}}\varphi\right)  _{k_{1}}\right\rangle
\left(  \left[  \mathbf{\Lambda}^{\left(  2\right)  }\left(  \Delta_{\varphi
}\right)  \right]  ^{-\frac{1}{2}}\mathbf{g}_{2\varphi}\right. \\
&  \left.  \otimes\mathbf{g}_{2\varphi}^{\dagger}\left[  \mathbf{\Lambda
}^{\left(  2\right)  }\left(  \Delta_{\varphi}\right)  \right]  ^{-\frac{1}%
{2}}\right)  _{j_{1}k_{1}l_{2}m_{2}}\left\langle \left(  \Delta_{\varphi
}^{-\frac{1}{2}}\varphi\right)  _{l_{2}}\left(  \Delta_{\varphi}^{-\frac{1}%
{2}}\varphi\right)  _{m_{2}}\right\vert
\end{align}


The contraction operator is defined by%
\begin{align}
L_{2}^{1}\left(  \left\vert w\wedge x\right\rangle \left\langle y\wedge
z\right\vert \right)   &  =\frac{1}{4}\left\{  \left\vert w\otimes x-x\otimes
w\right\rangle \left\langle y\otimes z-z\otimes y\right\vert \right\}
\nonumber\\
&  =\frac{1}{4}\left\{  \left\vert w\right\rangle \left\langle y\right\vert
\left\langle \left.  z\right\vert x\right\rangle -\left\vert w\right\rangle
\left\langle z\right\vert \left\langle \left.  y\right\vert x\right\rangle
-\left\vert x\right\rangle \left\langle y\right\vert \left\langle \left.
z\right\vert w\right\rangle +\left\vert x\right\rangle \left\langle
z\right\vert \left\langle \left.  y\right\vert w\right\rangle \right\}
\end{align}
In particular%
\begin{equation}
L_{2}^{1}\left(  \left\vert w\wedge x\right\rangle \left\langle w\wedge
x\right\vert \right)  =\frac{1}{4}\left\{  \left\vert w\right\rangle
\left\langle w\right\vert \left\langle \left.  x\right\vert x\right\rangle
-\left\vert w\right\rangle \left\langle x\right\vert \left\langle \left.
x\right\vert w\right\rangle -\left\vert x\right\rangle \left\langle
w\right\vert \left\langle \left.  w\right\vert x\right\rangle +\left\vert
x\right\rangle \left\langle x\right\vert \left\langle \left.  w\right\vert
w\right\rangle \right\}
\end{equation}
and if $\left\vert w\right\rangle $ and $\left\vert x\right\rangle $ are part
of an on basis then
\begin{equation}
L_{2}^{1}\left(  \left\vert w\wedge x\right\rangle \left\langle w\wedge
x\right\vert \right)  =\frac{1}{2}\left\{  \left\vert w\right\rangle
\left\langle w\right\vert +\left\vert x\right\rangle \left\langle x\right\vert
\right\}  .
\end{equation}%
\begin{align}
&  L_{2}^{1}\left(  \left\vert g\right\rangle \left\langle g\right\vert
\right)  =L_{2}^{1}\left(  \frac{1}{4}\sum_{\substack{1\leq j,k\leq r\\1\leq
l,m\leq r}}g_{jk}\bar{g}_{lm}\left\vert \varphi_{j}\varphi_{k}\right\rangle
\left\langle \varphi_{l}\varphi_{m}\right\vert \right) \\
&  =\frac{1}{16}\sum_{\substack{1\leq j,k\leq r\\1\leq l,m\leq r}}g_{jk}%
\bar{g}_{lm}\left\{  \left\vert \varphi_{j}\right\rangle \left\langle
\varphi_{l}\right\vert \delta_{km}-\left\vert \varphi_{j}\right\rangle
\left\langle \varphi_{m}\right\vert \delta_{kl}-\left\vert \varphi
_{k}\right\rangle \left\langle \varphi_{l}\right\vert \delta_{jm}+\left\vert
\varphi_{k}\right\rangle \left\langle \varphi_{m}\right\vert \delta
_{jl}\right\} \\
&  =\frac{1}{16}\sum_{\substack{1\leq j,k\leq r\\1\leq l,m\leq r}}\left\{
\left\vert \varphi_{j}\right\rangle \left\langle \varphi_{l}\right\vert
g_{jk}\bar{g}_{lk}+\left\vert \varphi_{j}\right\rangle \left\langle
\varphi_{m}\right\vert g_{jk}\bar{g}_{mk}+\left\vert \varphi_{k}\right\rangle
\left\langle \varphi_{l}\right\vert g_{jk}\bar{g}_{jl}+\left\vert \varphi
_{k}\right\rangle \left\langle \varphi_{m}\right\vert g_{jk}\bar{g}%
_{jm}\right\} \\
&  =\frac{1}{4}\sum_{\substack{1\leq j,k\leq r\\1\leq l,m\leq r}}g_{jk}\bar
{g}_{lm}\left\vert \varphi_{j}\right\rangle \left\langle \varphi
_{l}\right\vert =\frac{1}{4}\sum_{1\leq j,l\leq r}L_{2}^{1}\left(  \left\vert
g\right\rangle \left\langle g\right\vert \right)  _{jl}\left\vert \varphi
_{j}\right\rangle \left\langle \varphi_{l}\right\vert \\
&  L_{2}^{1}\left(  \left\vert g\right\rangle \left\langle g\right\vert
\right)  _{jl}=\frac{1}{4}\sum_{1\leq k\leq r}g_{jk}\bar{g}_{lk}%
\end{align}
%

\begin{align}
L_{2}^{1}\left(  \left\vert g\right\rangle \left\langle g\right\vert \right)
_{\varphi j_{1}l_{2}}  &  =\sum_{\substack{1\leq j_{1},k_{1}\leq r\\1\leq
l_{2},m_{2}\leq r}}\left\vert \left(  \Delta_{\varphi}^{-\frac{1}{2}}%
\varphi\right)  _{j_{1}}\right\rangle \left(  \left[  \mathbf{\Lambda
}^{\left(  2\right)  }\left(  \Delta_{\varphi}\right)  \right]  ^{-\frac{1}%
{2}}\mathbf{g}_{2\varphi}\otimes\mathbf{g}_{2\varphi}^{\dagger}\left[
\mathbf{\Lambda}^{\left(  2\right)  }\left(  \Delta_{\varphi}\right)  \right]
^{-\frac{1}{2}}\right)  _{j_{1}k_{1}l_{2}m_{2}}\\
&  \left\langle \left(  \Delta_{\varphi}^{-\frac{1}{2}}\varphi\right)
_{l_{2}}\right\vert \delta_{k_{1}m_{2}}%
\end{align}


\subsubsection{Antisymmetric operator}

This expansion of $\left\vert g\right\rangle $ defines an antisymmetric
operator in the orthonormalized basis
\begin{align}
\Upsilon_{\tilde{\varphi}}\left(  g\right)   &  =\frac{1}{2}\sum_{1\leq
j,k\leq r}\mathcal{R}_{2\tilde{\varphi}}\left(  g\right)  _{jk}\left\vert
\tilde{\varphi}_{j}\right\rangle \left\langle \tilde{\varphi}_{k}\right\vert
\equiv\frac{1}{2}\sum_{1\leq j,k\leq r}\mathcal{R}_{\tilde{\varphi}}\left(
g\right)  _{jk}\left\vert \tilde{\varphi}_{j}\right\rangle \left\langle
\tilde{\varphi}_{k}\right\vert \nonumber\\
&  =\frac{1}{2}\sum_{1\leq j,k\leq r}\mathbf{g}_{\tilde{\varphi}jk}\left\vert
\tilde{\varphi}_{j}\right\rangle \left\langle \tilde{\varphi}_{k}\right\vert ,
\end{align}
that can be expressed in the general basis as%
\begin{align}
&  \Upsilon_{\tilde{\varphi}}\left(  g\right)  =\frac{1}{2}\sum_{1\leq j,k\leq
r}\left(  \Delta_{\varphi}^{-\frac{1}{2}}\left(  \left[  \mathbf{\Lambda
}^{\left(  2\right)  }\left(  \Delta_{\varphi}\right)  \right]  ^{\frac{1}{2}%
}\mathcal{R}_{2\varphi}\left(  g\right)  \right)  \Delta_{\varphi}^{-\frac
{1}{2}}\right)  _{jk}\left\vert \varphi_{j}\right\rangle \left\langle
\varphi_{k}\right\vert \nonumber\\
&  \equiv\frac{1}{2}\sum_{1\leq j,k\leq r}\left(  \Delta_{\varphi}^{-\frac
{1}{2}}\mathcal{R}_{\tilde{\varphi}}\left(  g\right)  \Delta_{\varphi}%
^{-\frac{1}{2}}\right)  _{jk}\left\vert \varphi_{j}\right\rangle \left\langle
\varphi_{k}\right\vert \nonumber\\
&  =\frac{1}{2}\sum_{1\leq j,k\leq r}\mathbf{g}_{\tilde{\varphi}jk}\left\vert
\tilde{\varphi}_{j}\right\rangle \left\langle \tilde{\varphi}_{k}\right\vert
=\frac{1}{2}\sum_{1\leq j,k\leq r}\left(  \Delta_{\varphi}^{-\frac{1}{2}%
}\left[  \mathbf{\Lambda}^{\left(  2\right)  }\left(  \Delta_{\varphi}\right)
\right]  ^{-\frac{1}{2}}\mathbf{g}_{2\varphi}\Delta_{\varphi}^{-\frac{1}{2}%
}\right)  _{jk}\left\vert \varphi_{j}\right\rangle \left\langle \varphi
_{k}\right\vert ,
\end{align}
$\lceil$which is consistent with the relationships%
\begin{align*}
\mathcal{R}_{2\tilde{\varphi}}\left(  g\right)   &  =\left[  \mathbf{\Lambda
}^{\left(  2\right)  }\left(  \Delta_{\varphi}\right)  \right]  ^{-\frac{1}%
{2}}\mathcal{R}_{2\varphi}\left(  g\right) \\
\mathcal{R}_{2\varphi}\left(  g\right)   &  =\left[  \mathbf{\Lambda}^{\left(
2\right)  }\left(  \Delta_{\varphi}\right)  \right]  ^{-1}\mathbf{g}%
_{2\varphi}\\
\mathbf{g}_{2\tilde{\varphi}}  &  =\left[  \mathbf{\Lambda}^{\left(  2\right)
}\left(  \Delta_{\varphi}\right)  \right]  ^{-\frac{1}{2}}\mathbf{g}%
_{2\varphi}.
\end{align*}
$\lceil$%
\begin{align}
\mathbf{g}_{\varphi}  &  =\Delta_{\varphi}\mathcal{R}_{\varphi}\left(
\Upsilon_{\varphi}\left(  g\right)  \right)  \equiv\Delta_{\varphi}%
\mathcal{R}_{\varphi}\left(  g\right) \nonumber\\
\mathbf{g}_{2\varphi}  &  =\mathbf{\Lambda}^{\left(  2\right)  }\left(
\Delta_{\varphi}\right)  \mathcal{R}_{2\varphi}\left(  g\right) \nonumber\\
\mathcal{R}_{\varphi}\left(  g\right)   &  =\Delta_{\varphi}^{-\frac{1}{2}%
}\mathcal{R}_{\tilde{\varphi}}\left(  g\right)  \Delta_{\varphi}^{-\frac{1}%
{2}}=\Delta_{\varphi}^{-\frac{1}{2}}\left[  \mathbf{\Lambda}^{\left(
2\right)  }\left(  \Delta_{\varphi}\right)  \right]  ^{-\frac{1}{2}}%
\mathbf{g}_{2\varphi}\Delta_{\varphi}^{-\frac{1}{2}}\\
&  =\Delta_{\varphi}^{-\frac{1}{2}}\left[  \mathbf{\Lambda}^{\left(  2\right)
}\left(  \Delta_{\varphi}\right)  \right]  ^{\frac{1}{2}}\mathcal{R}%
_{2\varphi}\left(  g\right)  \Delta_{\varphi}^{-\frac{1}{2}}\\
\Upsilon_{\varphi}\left(  g\right)   &  =\left(  \left\vert \varphi
_{1}\right\rangle \left\langle \varphi_{1}\right\vert \cdots\left\vert
\varphi_{r}\right\rangle \left\langle \varphi_{r}\right\vert \right)
\mathcal{R}_{\varphi}\left(  g\right) \nonumber\\
\left\vert g\right\rangle  &  =\left(  \left\vert \varphi_{1}\varphi
_{1}\right\rangle \cdots\left\vert \varphi_{r}\varphi_{r}\right\rangle
\right)  \mathcal{R}_{\varphi}\left(  \Upsilon_{\varphi}\left(  g\right)
\right)  \rfloor
\end{align}


where
\begin{equation}
\mathbf{g}_{\varphi}=\left(
\begin{array}
[c]{c}%
\left\langle \varphi_{1}\right\vert \\
\vdots\\
\left\langle \varphi_{r}\right\vert
\end{array}
\right)  \Upsilon_{\varphi}\left(  g\right)  \left(  \left\vert \varphi
_{1}\right\rangle ,\cdots,\left\vert \varphi_{r}\right\rangle \right)
\end{equation}
and the metric matrix
\begin{equation}
\Delta_{\varphi}=\left(
\begin{array}
[c]{c}%
\left\langle \varphi_{1}\right\vert \\
\vdots\\
\left\langle \varphi_{r}\right\vert
\end{array}
\right)  \left(  \left\vert \varphi_{1}\right\rangle ,\cdots,\left\vert
\varphi_{r}\right\rangle \right)  .
\end{equation}
The norm of $\left\vert g\right\rangle $ is
\begin{align}
\left\Vert g\right\Vert _{\mathcal{H}^{2}}  &  =\left\langle \left.
g\right\vert g\right\rangle ^{\frac{1}{2}}=\frac{1}{2}\left(  \sum
_{\substack{1\leq j,k\leq s\\1\leq l,m\leq s}}\overline{\mathcal{R}_{2\varphi
}\left(  g\right)  }_{jk}\mathcal{R}_{2\varphi}\left(  g\right)
_{lm}\left\langle \left.  \varphi_{j}\wedge\varphi_{k}\right\vert \varphi
_{l}\wedge\varphi_{m}\right\rangle \right)  ^{\frac{1}{2}}\\
&  =\frac{1}{2\sqrt{2}}\left(  \sum_{\substack{1\leq j,k\leq s\\1\leq l,m\leq
s}}\overline{\mathcal{R}_{2\varphi}\left(  g\right)  }_{jk}\mathbf{\Lambda
}^{\left(  2\right)  }\left(  \Delta_{\varphi}\right)  _{jklm}\mathcal{R}%
_{2\varphi}\left(  g\right)  _{lm}\right)  ^{\frac{1}{2}},
\end{align}
while the HS norm of $\Upsilon_{\varphi}\left(  g\right)  $ is
\begin{align}
\left\Vert \Upsilon_{\varphi}\left(  g\right)  \right\Vert _{\mathcal{B}%
_{2}\left(  \mathcal{H}^{1}\right)  }  &  =\left(  \operatorname*{Tr}\left\{
\Upsilon_{\varphi}\left(  g\right)  ^{\dagger}\Upsilon_{\varphi}\left(
g\right)  \right\}  \right)  ^{\frac{1}{2}}\\
&  =\left\{  \frac{1}{4}\sum_{\substack{1\leq j,k\leq r\\1\leq l,m\leq
r}}\overline{\mathcal{R}_{\varphi}\left(  g\right)  _{jk}}\mathcal{R}%
_{\varphi}\left(  g\right)  _{lm}\left\langle \left.  \varphi_{j}\right\vert
\varphi_{l}\right\rangle \left\langle \left.  \varphi_{m}\right\vert
\varphi_{k}\right\rangle \right\}  ^{\frac{1}{2}}\\
&  =\frac{1}{2}\left\{  \sum_{\substack{1\leq j,k\leq r\\1\leq l,m\leq
r}}\mathcal{R}_{\varphi}\left(  g\right)  _{kj}^{\dagger}\Delta_{\varphi
jl}\mathcal{R}_{\varphi}\left(  g\right)  _{lm}\Delta_{\varphi mk}\right\}
^{\frac{1}{2}}\\
&  =\frac{1}{2}\left(  \operatorname*{Tr}\left\{  \mathcal{R}_{\varphi}\left(
g\right)  ^{\dagger}\Delta_{\varphi}\mathcal{R}_{\varphi}\left(  g\right)
\Delta_{\varphi}\right\}  \right)  ^{\frac{1}{2}}%
\end{align}%
\begin{align}
\mathcal{R}_{\varphi}\left(  g\right)   &  =\Delta_{\varphi}^{-1}%
\mathbf{g}_{\varphi},\\
\mathcal{R}_{\varphi}\left(  g\right)  ^{\dagger}  &  =\mathbf{g}_{\varphi
}^{\dagger}\Delta_{\varphi}^{-1}%
\end{align}%
\begin{equation}
\frac{1}{2}\left(  \operatorname*{Tr}\left\{  \mathbf{g}_{\varphi}^{\dagger
}\Delta_{\varphi}^{-1}\Delta_{\varphi}\Delta_{\varphi}^{-1}\mathbf{g}%
_{\varphi}\Delta_{\varphi}\right\}  \right)  ^{\frac{1}{2}}=\frac{1}{2}\left(
\operatorname*{Tr}\left\{  \mathbf{g}_{\varphi}^{\dagger}\Delta_{\varphi}%
^{-1}\mathbf{g}_{\varphi}\Delta_{\varphi}\right\}  \right)  ^{\frac{1}{2}}%
\end{equation}


The matrix, $\mathbb{G}_{\varphi},$ of the symplectic form wrt the conb
$\left\{  \left.  \left\vert \varphi_{j}\right\rangle \right\vert 1\leq
j\leq2s=r\right\}  $ of $\mathcal{H}^{1}$ is defined as
\begin{equation}
\mathbb{G}_{\varphi jk}=g\left(  \varphi_{j},\varphi_{k}\right)  =\left\langle
\left.  g\right\vert \varphi_{j}\wedge\varphi_{k}\right\rangle .
\end{equation}
The value of this form does not depend on the way it is expressed so wlog one
can express $g\left(  \varphi_{j},\varphi_{k}\right)  $ as
\begin{align}
g\left(  \varphi_{l},\varphi_{m}\right)   &  =\frac{1}{2}\sum_{1\leq j,k\leq
s}\mathcal{R}_{2\varphi}\left(  g\right)  _{jk}\left\langle \left.
\varphi_{j}\wedge\varphi_{k}\right\vert \varphi_{l}\wedge\varphi
_{m}\right\rangle \\
&  =\frac{1}{4}\sum_{1\leq j,k\leq s}\mathcal{R}_{2\varphi}\left(  g\right)
_{jk}\det\left(
\begin{array}
[c]{cc}%
\left\langle \left.  \varphi_{j}\right\vert \varphi_{l}\right\rangle  &
\left\langle \left.  \varphi_{j}\right\vert \varphi_{m}\right\rangle \\
\left\langle \left.  \varphi_{k}\right\vert \varphi_{l}\right\rangle  &
\left\langle \left.  \varphi_{k}\right\vert \varphi_{m}\right\rangle
\end{array}
\right) \\
&  =\frac{1}{4}\sum_{1\leq j,k\leq s}\mathcal{R}_{2\varphi}\left(  g\right)
_{jk}\det\left(
\begin{array}
[c]{cc}%
\Delta_{\varphi jl} & \Delta_{\varphi jm}\\
\Delta_{\varphi kl} & \Delta_{\varphi km}%
\end{array}
\right) \\
&  =\frac{1}{4}\sum_{1\leq j,k\leq s}\mathcal{R}_{2\varphi}\left(  g\right)
_{jk}\mathbf{\Lambda}^{\left(  2\right)  }\left(  \Delta_{\varphi}\right)
_{jklm}%
\end{align}
$\lceil$Note that
\begin{align}
u\wedge v  &  =\frac{1}{2}\left(  u\otimes v-v\otimes u\right) \\
\left\langle \left.  a\wedge b\right\vert u\wedge v\right\rangle  &  =\frac
{1}{4}\left\{  \left\langle \left.  a\otimes b\right\vert u\otimes
v\right\rangle -\left\langle \left.  a\otimes b\right\vert v\otimes
u\right\rangle \right. \\
&  \left.  -\left\langle \left.  b\otimes a\right\vert u\otimes v\right\rangle
+\left\langle \left.  b\otimes a\right\vert v\otimes u\right\rangle \right\}
\\
&  =\frac{1}{2}\left\{  \left\langle \left.  a\right\vert u\right\rangle
\left\langle \left.  b\right\vert v\right\rangle -\left\langle \left.
a\right\vert v\right\rangle \left\langle \left.  b\right\vert u\right\rangle
\right\} \\
&  =\frac{1}{2}\det\left\{
\begin{array}
[c]{cc}%
\left\langle \left.  a\right\vert u\right\rangle  & \left\langle \left.
a\right\vert v\right\rangle \\
\left\langle \left.  b\right\vert u\right\rangle  & \left\langle \left.
b\right\vert v\right\rangle
\end{array}
\right\}  \rfloor
\end{align}
This should be compared to
\begin{equation}
\left\vert g\right\rangle =\left(  \left\vert \varphi_{1}\wedge\varphi
_{2}\right\rangle ,\cdots,\left\vert \varphi_{r-1}\wedge\varphi_{r}%
\right\rangle \right)  \mathcal{R}_{2\varphi}\left(  g\right)
\end{equation}%
\begin{align}
\mathbf{g}_{2\varphi}  &  =\left\langle \left.  \left(
\begin{array}
[c]{c}%
\left\langle \varphi_{1}\wedge\varphi_{2}\right\vert \\
\vdots\\
\left\langle \varphi_{r-1}\wedge\varphi_{r}\right\vert
\end{array}
\right)  \right\vert g\right\rangle \\
&  =\left\langle \left.  \left(
\begin{array}
[c]{c}%
\left\langle \varphi_{1}\wedge\varphi_{2}\right\vert \\
\vdots\\
\left\langle \varphi_{r-1}\wedge\varphi_{r}\right\vert
\end{array}
\right)  \right\vert \left(  \left\vert \varphi_{1}\wedge\varphi
_{2}\right\rangle ,\cdots,\left\vert \varphi_{r-1}\wedge\varphi_{r}%
\right\rangle \right)  \right\rangle \mathcal{R}_{2\varphi}\left(  g\right) \\
&  =\mathbf{\Lambda}^{\left(  2\right)  }\left(  \Delta_{\varphi}\right)
\mathcal{R}_{2\varphi}\left(  g\right)  ,
\end{align}


As discussed in the following the transformation properties of $\mathcal{R}%
_{\varphi}\left(  \Lambda\left(  g\right)  \right)  ,$ $\mathbf{g}_{\varphi},$
$\mathbf{g}_{2\varphi},\mathbb{G}_{\varphi},$ and $\mathcal{R}_{2\varphi
}\left(  g\right)  $ are not the same wrt a change of basis of $\mathcal{H}%
^{1},$ where
\begin{align}
\Lambda\left(  g\right)   &  =\Lambda\left(  g\right)  \left(  \left\vert
\varphi_{1}\right\rangle ,\cdots,\left\vert \varphi_{r}\right\rangle \right)
\Delta_{\varphi}^{-1}\left(
\begin{array}
[c]{c}%
\left\langle \varphi_{1}\right\vert \\
\vdots\\
\left\langle \varphi_{r}\right\vert
\end{array}
\right) \\
&  =\left(  \left\vert \varphi_{1}\right\rangle ,\cdots,\left\vert \varphi
_{r}\right\rangle \right)  \mathcal{R}_{\varphi}\left(  \Lambda\left(
g\right)  \right)  \Delta_{\varphi}^{-1}\left(
\begin{array}
[c]{c}%
\left\langle \varphi_{1}\right\vert \\
\vdots\\
\left\langle \varphi_{r}\right\vert
\end{array}
\right)  ,\\
\mathbf{g}_{\varphi}  &  =\mathcal{R}_{\varphi}\left(  \Lambda\left(
g\right)  \right)  \Delta_{\varphi}^{-1}=\left(
\begin{array}
[c]{c}%
\left\langle \varphi_{1}\right\vert \\
\vdots\\
\left\langle \varphi_{r}\right\vert
\end{array}
\right)  \Lambda\left(  g\right)  \left(  \left\vert \varphi_{1}\right\rangle
,\cdots,\left\vert \varphi_{r}\right\rangle \right)  ,
\end{align}
The symplectic group $Sp\left(  2s,\mathbb{C}\right)  $ has a representation
on $\mathcal{H}^{1}$ given by
\begin{equation}
Sp\left(  2s,\mathbb{C}\right)  _{\mathcal{H}^{1}}=\left\{  \left.
T\right\vert g\left(  T\varphi_{1},T\varphi_{2}\right)  =g\left(  \varphi
_{1},\varphi_{2}\right)  \forall\left\vert \varphi_{1}\right\rangle
,\left\vert \varphi_{2}\right\rangle \in\mathcal{H}^{1}\right\}
\end{equation}
which induces the representations $\mathcal{R}_{\varphi}\left(  Sp\left(
2s,\mathbb{C}\right)  \right)  $ (the basis $\left\{  \left.  \left\vert
\varphi_{j}\right\rangle \right\vert 1\leq j\leq2s=r\right\}  $ does not have
to be orthonormal).

\subsection{Antisymmetric Operators and Geminals}

The space of antisymmetric Hilbert-Schmidt operators, $\mathcal{B}_{2a}\left(
\mathcal{H}^{1}\right)  $ can be defined as
\begin{align}
\mathcal{B}_{2a}\left(  \mathcal{H}^{1}\right)   &  =\left\{  \left.
A=\sum_{1\leq j,k\leq r}A_{\psi jk}\left\vert \psi_{j}\right\rangle
\left\langle \psi_{k}\right\vert \right\vert A_{\psi jk}=-A_{\psi kj},1\leq
j,k\leq r,A_{\psi jk}\in\mathbb{C},\right. \nonumber\\
&  \left.  \left\{  \left\vert \left\vert \psi_{j}\right\rangle \right.  1\leq
j\leq r\right\}  \text{ a fixed conb for }\mathcal{H}^{1}\right\} \nonumber\\
&  =\left\{  \mathbb{T}_{\psi}\mathbb{A}_{\psi}\mathbb{T}_{\psi}^{t}\left\vert
T\in\frac{GL\left(  2s,\mathbb{C}\right)  _{\mathcal{H}^{1}}}{Sp\left(
2s,\mathbb{C}\right)  _{\mathcal{H}^{1}}}\right.  \right\}  .
\end{align}


\subsubsection{Change of Basis of Operators}

Let $\left\{  \left\vert \left\vert \chi_{j}\right\rangle \right.  1\leq j\leq
r\right\}  $ be a basis of $\mathcal{H}^{1}$ that is not in general normalized
nor orthogonal then it produces the following resolution of identity
\begin{equation}
I=\sum_{1\leq j,k\leq r}\left\vert \chi_{j}\right\rangle \Delta_{\chi jk}%
^{-1}\left\langle \chi_{k}\right\vert
\end{equation}
$\lceil$ This can be seen from%
\begin{equation}
\sum_{1\leq j,k\leq r}\left\vert \chi_{j}\right\rangle \Delta_{\chi jk}%
^{-1}\left\langle \chi_{k}\right\vert \left\vert \chi_{m}\right\rangle
=\sum_{1\leq j,k\leq r}\left\vert \chi_{j}\right\rangle \Delta_{\chi jk}%
^{-1}\Delta_{\chi km}=\sum_{1\leq j\leq r}\left\vert \chi_{j}\right\rangle
\delta_{jm}=\left\vert \chi_{m}\right\rangle
\end{equation}
$\rfloor$

leading to
\begin{equation}
I=\left(  \left\vert \chi_{1}\right\rangle ,\cdots,\left\vert \chi
_{r}\right\rangle \right)  \Delta_{\chi}^{-1}\left(
\begin{array}
[c]{c}%
\left\langle \chi_{1}\right\vert \\
\vdots\\
\left\langle \chi_{r}\right\vert
\end{array}
\right)  .
\end{equation}
The operator $\Lambda\left(  g\right)  $ can then be expanded as
\begin{equation}
\Lambda\left(  g\right)  =\Lambda\left(  g\right)  \left(  \left\vert \chi
_{1}\right\rangle ,\cdots,\left\vert \chi_{r}\right\rangle \right)
\Delta_{\chi}^{-1}\left(
\begin{array}
[c]{c}%
\left\langle \chi_{1}\right\vert \\
\vdots\\
\left\langle \chi_{r}\right\vert
\end{array}
\right)  =\left(  \left\vert \chi_{1}\right\rangle ,\cdots,\left\vert \chi
_{r}\right\rangle \right)  \mathcal{R}_{\chi}\left(  g\right)  \Delta_{\chi
}^{-1}\left(
\begin{array}
[c]{c}%
\left\langle \chi_{1}\right\vert \\
\vdots\\
\left\langle \chi_{r}\right\vert
\end{array}
\right)  ,
\end{equation}
i.e.
\begin{equation}
\Lambda\left(  g\right)  =\sum_{1\leq j,k\leq r}\left\vert \chi_{j}%
\right\rangle \left(  \mathcal{R}_{\chi}\left(  g\right)  \Delta_{\chi}%
^{-1}\right)  _{jk}\left\langle \chi_{k}\right\vert .
\end{equation}
The above expansion gives that
\begin{equation}
\mathbf{g}_{\chi}=\Delta_{\chi}\mathcal{R}_{\chi}\left(  g\right)
\equiv\mathcal{R}_{\chi}\left(  g\right)  =\Delta_{\chi}^{-1}\mathbf{g}_{\chi
},
\end{equation}
where the matrix of $g$ wrt the basis $\left\vert \mathbf{\chi}\right\rangle $
is given as
\begin{equation}
\mathbf{g}_{\chi}=\left(
\begin{array}
[c]{c}%
\left\langle \chi_{1}\right\vert \\
\vdots\\
\left\langle \chi_{r}\right\vert
\end{array}
\right)  g\left(  \left\vert \chi_{1}\right\rangle ,\cdots,\left\vert \chi
_{r}\right\rangle \right)  .
\end{equation}
This leads to
\begin{equation}
\Delta_{\chi}^{-1}\mathbf{g}_{\chi}\Delta_{\chi}^{-1}=\mathcal{R}_{\chi
}\left(  g\right)  \Delta_{\chi}^{-1}%
\end{equation}
and
\begin{equation}
\Lambda\left(  g\right)  =\sum_{1\leq j,k\leq r}\left\vert \chi_{j}%
\right\rangle \left(  \Delta_{\chi}^{-1}\mathbf{g}_{\chi}\Delta_{\chi}%
^{-1}\right)  _{jk}\left\langle \chi_{k}\right\vert .
\end{equation}


$\lceil$%
\begin{equation}
\mathbf{g}_{\chi}=\left(
\begin{array}
[c]{c}%
\left\langle \chi_{1}\right\vert \\
\vdots\\
\left\langle \chi_{r}\right\vert
\end{array}
\right)  g\left(  \left\vert \chi_{1}\right\rangle ,\cdots,\left\vert \chi
_{r}\right\rangle \right)  =\Delta_{\chi}\Delta_{\chi}^{-1}\mathbf{g}_{\chi
}\Delta_{\chi}^{-1}\Delta_{\chi}=\mathbf{g}_{\chi}.
\end{equation}
$\rfloor$

A change of basis in $\mathcal{H}^{1}$ produces the following effect on the
matrices $\Delta_{\chi}$ and $\mathbf{g}_{\chi}$
\begin{align}
\Delta_{\chi}  &  =\mathbb{S}_{\chi\psi}^{\dagger}\Delta_{\psi}\mathbb{S}%
_{\chi\psi}\nonumber\\
\mathbf{g}_{\chi}  &  =\mathbb{S}_{\chi\psi}^{\dagger}\mathbf{g}_{\psi
}\mathbb{S}_{\chi\psi},
\end{align}
as can be seen from
\begin{align}
\Delta_{\chi}  &  =\mathbb{S}_{\chi\psi}^{\dagger}\left(
\begin{array}
[c]{c}%
\left\langle \psi_{1}\right\vert \\
\vdots\\
\left\langle \psi_{r}\right\vert
\end{array}
\right)  \left(  \left\vert \psi_{1}\right\rangle ,\cdots,\left\vert \psi
_{r}\right\rangle \right)  \mathbb{S}_{\chi\psi}\nonumber\\
&  =\mathbb{S}_{\chi\psi}^{\dagger}\Delta_{\psi}\mathbb{S}_{\chi\psi}.
\end{align}
This is consistent with
\begin{align}
\Lambda_{\chi}\left(  g\right)   &  =\sum_{1\leq j,k\leq r}\left\vert \chi
_{j}\right\rangle \left(  \Delta_{\chi}^{-1}\mathbf{g}_{\chi}\Delta_{\chi
}^{-1}\right)  _{jk}\left\langle \chi_{k}\right\vert \nonumber\\
&  =\sum_{1\leq j,k\leq r}\left\vert S_{\chi\psi}\psi_{j}\right\rangle \left(
\mathbb{S}_{\chi\psi}^{-1}\Delta_{\psi}^{-1}\mathbb{S}_{\chi\psi}^{-\dagger
}\mathbb{S}_{\chi\psi}^{\dagger}\mathbf{g}_{\psi}\mathbb{S}_{\chi\psi
}\mathbb{S}_{\chi\psi}^{-1}\Delta_{\psi}^{-1}\mathbb{S}_{\chi\psi}^{-\dagger
}\right)  _{jk}\left\langle S_{\chi\psi}\psi_{k}\right\vert \nonumber\\
&  =\sum_{1\leq j,k\leq r}\left\vert \psi_{j}\right\rangle \left(
\mathbb{S}_{\chi\psi}\mathbb{S}_{\chi\psi}^{-1}\Delta_{\psi}^{-1}%
\mathbb{S}_{\chi\psi}^{-\dagger}\mathbb{S}_{\chi\psi}^{\dagger}\mathbf{g}%
_{\psi}\mathbb{S}_{\chi\psi}\mathbb{S}_{\chi\psi}\Delta_{\psi}^{-1}%
\mathbb{S}_{\chi\psi}^{-\dagger}\mathbb{S}_{\chi\psi}^{\dagger}\right)
_{jk}\left\langle \psi_{k}\right\vert \nonumber\\
&  =\sum_{1\leq j,k\leq r}\left\vert \psi_{j}\right\rangle \left(
\Delta_{\psi}^{-1}\mathbf{g}_{\psi}\Delta_{\psi}^{-1}\right)  _{jk}%
\left\langle \psi_{k}\right\vert .
\end{align}
As
\begin{equation}
\mathbf{g}_{\chi}=\mathbb{S}_{\chi\psi}^{\dagger}\mathbf{g}_{\psi}%
\mathbb{S}_{\chi\psi}%
\end{equation}
the antisymmetry
\begin{equation}
\mathbf{g}_{\chi}=-\mathbf{g}_{\chi}^{t}\text{ }%
\end{equation}
is conserved if and only if
\begin{equation}
\mathbb{S}_{\chi\psi}=\mathbb{\bar{S}}_{\chi\psi}%
\end{equation}
i.e. $\mathbb{S}_{\chi\psi}$ is a real matrix.

In the following we will reorder the canonical basis so that the
representation of $g$ is given by
\begin{equation}
\Lambda\left(  g\right)  \left(  \left\vert \psi_{1}\right\rangle
,\cdots,\left\vert \psi_{s}\right\rangle ,\left\vert \psi_{-1}\right\rangle
,\cdots,\left\vert \psi_{-s}\right\rangle \right)  =\left(  \left\vert
\psi_{1}\right\rangle ,\cdots,\left\vert \psi_{s}\right\rangle ,\left\vert
\psi_{-1}\right\rangle ,\cdots,\left\vert \psi_{-s}\right\rangle \right)
\mathcal{R}_{\psi}\left(  g\right)  ,
\end{equation}
i.e.
\begin{equation}
\Lambda\left(  g\right)  \left\vert \psi_{j}\right\rangle =\sum_{s\geq
k\geq-s}\left\vert \psi_{k}\right\rangle \mathcal{R}_{\psi}\left(  g\right)
_{kj},
\end{equation}
thus
\begin{equation}
\mathbf{g}_{\psi}=\Delta_{\psi}\mathcal{R}_{\psi}\left(  g\right)
\end{equation}
and
\begin{equation}
\mathcal{R}_{\psi}\left(  g\right)  =\Delta_{\psi}^{-1}\mathbf{g}_{\psi},
\end{equation}
where the matrix of $g$ is given by
\begin{equation}
\mathbf{g}_{\psi}=\left(
\begin{array}
[c]{c}%
\left\langle \psi_{1}\right\vert \\
\vdots\\
\left\langle \psi_{s}\right\vert \\
\left\langle \psi_{-1}\right\vert \\
\vdots\\
\left\langle \psi_{-s}\right\vert
\end{array}
\right)  g\left(  \left\vert \psi_{1}\right\rangle ,\cdots,\left\vert \psi
_{s}\right\rangle ,\left\vert \psi_{-1}\right\rangle ,\cdots,\left\vert
\psi_{-s}\right\rangle \right)
\end{equation}
and the metric matrix
\begin{equation}
\Delta_{\psi}=\left(
\begin{array}
[c]{c}%
\left\langle \psi_{1}\right\vert \\
\vdots\\
\left\langle \psi_{s}\right\vert \\
\left\langle \psi_{-1}\right\vert \\
\vdots\\
\left\langle \psi_{-s}\right\vert
\end{array}
\right)  \left(  \left\vert \psi_{1}\right\rangle ,\cdots,\left\vert \psi
_{s}\right\rangle ,\left\vert \psi_{-1}\right\rangle ,\cdots,\left\vert
\psi_{-s}\right\rangle \right)  .
\end{equation}


In the cannonical representation the preceding leads to
\begin{align}
\mathcal{R}_{\psi}\left(  g\right)   &  =\left(
\begin{array}
[c]{cc}%
\mathbb{O} & \mathbb{C}\\
-\mathbb{C} & \mathbb{O}%
\end{array}
\right) \\
\Delta_{\psi}  &  =\left(
\begin{array}
[c]{cc}%
\mathbb{I} & \mathbb{O}\\
\mathbb{O} & \mathbb{I}%
\end{array}
\right) \\
\mathbf{g}_{\psi}  &  =\left(
\begin{array}
[c]{cc}%
\mathbb{O} & \mathbb{C}\\
-\mathbb{C} & \mathbb{O}%
\end{array}
\right)
\end{align}
and
\begin{align}
\mathcal{R}_{\eta}\left(  g\right)   &  =\left(
\begin{array}
[c]{cc}%
\mathbb{O} & \mathbb{I}\\
-\mathbb{I} & \mathbb{O}%
\end{array}
\right) \\
\Delta_{\eta}  &  =\left(
\begin{array}
[c]{cc}%
\mathbb{C} & \mathbb{O}\\
\mathbb{O} & \mathbb{C}%
\end{array}
\right) \\
\mathbf{g}_{\eta}  &  =\left(
\begin{array}
[c]{cc}%
\mathbb{O} & \mathbb{C}\\
-\mathbb{C} & \mathbb{O}%
\end{array}
\right)  .
\end{align}


\subsubsection{Change of Basis of Geminals}

The matrix $\mathbb{G}_{\varphi}=$ $\left\{  \left.  g\left(  \varphi
_{i},\varphi_{j}\right)  \right\vert -1\leq i,j\leq s\right\}  $ of the
antisymmetric form $g$ wrt to the conb $\left\{  \left.  \varphi
_{j}\right\vert -1\leq j\leq s\right\}  $ given by
\begin{equation}
g\left(  \varphi_{i},\varphi_{j}\right)  =\left\langle \left.  g\right\vert
\varphi_{i}\wedge\varphi_{j}\right\rangle .
\end{equation}
The effect produced in $\mathcal{H}^{2}$ by a change of basis in
$\mathcal{H}^{1}$ on this matrix is
\begin{equation}
\mathbb{G}_{\varphi}=\mathbb{S}_{\chi\psi}^{t}\mathbb{G}_{\chi}\mathbb{S}%
_{\chi\psi}%
\end{equation}
in contrast to effect in $\mathcal{B}_{2a}\left(  \mathcal{H}^{1}\right)  .$
The matrices $\mathbb{G}_{\psi},$ $\mathbb{G}_{\eta}$ in normalized and
unnormalized cannonical basis are
\begin{equation}
\mathbb{G}_{\psi}=\left(
\begin{array}
[c]{cc}%
\mathbb{O} & \mathbb{C}\\
-\mathbb{C} & \mathbb{O}%
\end{array}
\right)
\end{equation}
In general
\begin{equation}
g\left(  u,v\right)  =\left\langle \left.  g\right\vert u\wedge v\right\rangle
=\mathbf{u}_{\varphi}^{t}\mathbb{G}_{\varphi}\mathbf{v}_{\varphi}\mathbf{,}%
\end{equation}
where
\begin{align}
\left\vert v\right\rangle  &  =\left(  \left\vert \varphi_{1}\right\rangle
,\cdots,\left\vert \varphi_{s}\right\rangle ,\left\vert \varphi_{-1}%
\right\rangle ,\cdots,\left\vert \varphi_{-s}\right\rangle \right)
\mathbf{v}_{\varphi}\\
\left\vert u\right\rangle  &  =\left(  \left\vert \varphi_{1}\right\rangle
,\cdots,\left\vert \varphi_{s}\right\rangle ,\left\vert \varphi_{-1}%
\right\rangle ,\cdots,\left\vert \varphi_{-s}\right\rangle \right)
\mathbf{u}_{\varphi}.
\end{align}


\subsubsection{Mapping between the Spaces}

Consider the map $M_{a}$ between $\mathcal{H}^{2}$ and the space of
antisymmetric Hilbert-Schmidt operators, $\mathcal{B}_{2a}\left(
\mathcal{H}^{1}\right)  ,$ i.e.
\begin{equation}
M_{a}:\mathcal{H}^{2}\rightarrow\mathcal{B}_{2a}\left(  \mathcal{H}%
^{1}\right)
\end{equation}
given by%
\begin{equation}
M_{a}\left(  \left\vert g\right\rangle \right)  =M_{a}\left(  \sum_{1\leq
j\leq s}c_{j}\left\vert \psi_{j}\psi_{-j}\right\rangle \right)  =\sum_{1\leq
j\leq s}c_{j}\left\{  \left\vert \psi_{j}\right\rangle \left\langle \psi
_{-j}\right\vert -\left\vert \psi_{-j}\right\rangle \left\langle \psi
_{j}\right\vert \right\}  .
\end{equation}
One can then observe that
\begin{align}
&  M_{a}\left(  \sum_{1\leq j\leq s}c_{j}\left\vert \psi_{j}\psi
_{-j}\right\rangle \right)  =M_{a}\left(  \sum_{1\leq j\leq s}c_{j}\left\vert
e^{i\theta_{j}}\psi_{j}e^{-i\theta_{j}}\psi_{-j}\right\rangle \right)
\nonumber\\
&  =\sum_{1\leq j\leq s}c_{j}\left\{  e^{2i\theta_{j}}\left\vert \psi
_{j}\right\rangle \left\langle \psi_{-j}\right\vert -e^{-2i\theta_{j}%
}\left\vert \psi_{-j}\right\rangle \left\langle \psi_{j}\right\vert \right\}
\end{align}
thus $M_{a}$ is multivalued. This can also be expressed using
\begin{align}
\left\vert \chi_{j}\left(  \theta_{j}\right)  \right\rangle  &  =\left\vert
e^{i\theta_{j}}\psi_{j}\right\rangle \nonumber\\
\left\vert \chi_{-j}\left(  \theta_{j}\right)  \right\rangle  &  =\left\vert
e^{-i\theta_{j}}\psi_{-j}\right\rangle
\end{align}
producing the antisymmetric operator
\begin{equation}
\sum_{1\leq j\leq s}c_{j}\left\{  \left\vert \chi_{j}\left(  \theta
_{j1}\right)  \right\rangle \left\langle \chi_{-j}\left(  \theta_{j1}\right)
\right\vert -\left\vert \chi_{-j}\left(  \theta_{j1}\right)  \right\rangle
\left\langle \chi_{j}\left(  \theta_{j1}\right)  \right\vert \right\}  .
\end{equation}
This transformation corresponds to the one parameter subgroup $SU\left(
1,1\right)  $ of $SU(2)$ generated by
\begin{equation}
u\left(  \theta\right)  =e^{i\theta_{j}\left(  \left\vert \psi_{j}%
\right\rangle \left\langle \psi_{j}\right\vert -\left\vert \psi_{-j}%
\right\rangle \left\langle \psi_{-j}\right\vert \right)  }.
\end{equation}


If we consider the "inverse" map $M_{b}$ defined as
\begin{equation}
M_{b}\left(  \sum_{1\leq j\leq s}c_{j}\left\{  \left\vert \psi_{j}%
\right\rangle \left\langle \psi_{-j}\right\vert -\left\vert \psi
_{-j}\right\rangle \left\langle \psi_{j}\right\vert \right\}  \right)
=\sum_{1\leq j\leq s}c_{j}\left\vert \psi_{j}\psi_{-j}\right\rangle
=\left\vert g\right\rangle
\end{equation}
then
\begin{align}
&  M_{b}\left(  e^{i\theta_{j}\left(  \left\vert \psi_{j}\right\rangle
\left\langle \psi_{j}\right\vert -\left\vert \psi_{-j}\right\rangle
\left\langle \psi_{-j}\right\vert \right)  }\sum_{1\leq j\leq s}c_{j}\left\{
\left\vert \psi_{j}\right\rangle \left\langle \psi_{-j}\right\vert -\left\vert
\psi_{-j}\right\rangle \left\langle \psi_{j}\right\vert \right\}  \right)
\nonumber\\
&  =M_{b}\left(  \sum_{1\leq j\leq s}c_{j}\left\{  e^{2i\theta_{j}}\left\vert
\psi_{j}\right\rangle \left\langle \psi_{-j}\right\vert -e^{-2i\theta_{j}%
}\left\vert \psi_{-j}\right\rangle \left\langle \psi_{j}\right\vert \right\}
\right) \nonumber\\
&  =\sum_{1\leq j\leq s}c_{j}\left\vert \psi_{j}\psi_{-j}\right\rangle .
\end{align}
The phase factor can be absorbed into the GSO's as
\begin{align}
A  &  =\sum_{1\leq j\leq s}c_{j}\left\{  e^{2i\theta_{j}}\left\vert \psi
_{j}\right\rangle \left\langle \psi_{-j}\right\vert -e^{-2i\theta_{j}%
}\left\vert \psi_{-j}\right\rangle \left\langle \psi_{j}\right\vert \right\}
\nonumber\\
&  =\sum_{1\leq j\leq s}c_{j}\left\{  \left\vert e^{i\theta_{j}}\psi
_{j}\right\rangle \left\langle e^{-i\theta_{j}}\psi_{-j}\right\vert
-\left\vert e^{-i\theta_{j}}\psi_{-j}\right\rangle \left\langle e^{i\theta
_{j}}\psi_{j}\right\vert \right\}  .
\end{align}


In general
\begin{align}
M_{b}\left(  \mathcal{R}_{\mathcal{H}^{1}}\left(  u\left(  \theta\right)
\right)  A\right)   &  =M_{b}\left(  A\right)  ,\quad A\in\mathcal{B}%
_{2a}\left(  \mathcal{H}^{1}\right)  ,u\left(  \theta\right)  \in SU\left(
1,1\right)  ^{s}\nonumber\\
u\left(  \theta\right)   &  :\ker\left\{  M_{b}\right\}  \rightarrow
\ker\left\{  M_{b}\right\} \nonumber\\
\mathcal{R}_{\mathcal{H}^{2}}\left(  u\left(  \theta\right)  \right)
M_{b}\left(  A\right)   &  \neq M_{b}\left(  \mathcal{R}_{\mathcal{H}^{1}%
}\left(  u\left(  \theta\right)  \right)  A\right)  ,u\left(  \theta\right)
\in SU\left(  1,1\right)  ^{s},
\end{align}
but
\begin{equation}
\mathcal{R}_{\mathcal{H}^{2}}\left(  u\left(  \theta\right)  \right)
M_{b}\left(  A\right)  =M_{b}\left(  \mathcal{R}_{\mathcal{H}^{1}}\left(
u\left(  \theta\right)  \right)  A\right)  ,u\left(  \theta\right)  \in
SU\left(  2\right)  ^{s}%
\end{equation}
and $M_{b}$ is singular (it is surjective but not injective). This gives
\begin{equation}
M_{b}\left(  \sum_{1\leq j\leq s}c_{j}\left\{  \left(  e^{2i\theta_{1j}%
}-e^{2i\theta_{2j}}\right)  \left\vert \psi_{j}\right\rangle \left\langle
\psi_{-j}\right\vert -\left(  e^{-2i\theta_{1j}}-e^{-2i\theta_{2j}}\right)
\left\vert \psi_{-j}\right\rangle \left\langle \psi_{j}\right\vert \right\}
\right)  =0.
\end{equation}


\subsubsection{Invariance Groups}

One should note that if $\left\vert g\right\rangle $ is non degenerate the
maximal compact invariance group is the special unitary group $SU\left(
2\right)  ^{s},$ whose complexification is $SL\left(  2,\mathbb{C}\right)
^{s}$ i.e.
\begin{align}
&  \sum_{1\leq j\leq s}c_{j}\left\vert u\psi_{j}u\psi_{-j}\right\rangle
=\sum_{1\leq j\leq s}c_{j}\left\vert \psi_{j}\psi_{-j}\right\rangle ,u\in
SU\left(  2\right)  ^{s},\nonumber\\
&  \sum_{1\leq j\leq s}c_{j}\left\vert v\psi_{j}v\psi_{-j}\right\rangle
=\sum_{1\leq j\leq s}c_{j}\left\vert \psi_{j}\psi_{-j}\right\rangle ,v\in
SL\left(  2,\mathbb{C}\right)  ^{s},
\end{align}
while if the degeneracy is twofold then the maximal compact invariance group
is $SU\left(  2\right)  ^{s-2}\times USp\left(  4\right)  $, while the non
compact one is $SL\left(  2,\mathbb{C}\right)  ^{s-2}\times Sp\left(
4,\mathbb{C}\right)  $ etc until one has an $s$-fold degeneracy then the
maximal invariance are groups $USp\left(  2s\right)  $ and $Sp\left(
2s,\mathbb{C}\right)  .$

The invariance group for a non degenerate antisymmetric operator is $SU\left(
1,1\right)  ^{s},$ which changes with degeneracy from $SU\left(  1,1\right)
^{s-2}\times$ $SU\left(  2,2\right)  $ to $SU\left(  s,s\right)  .$ The
complexification of $SU\left(  1,1\right)  $ is $Sp\left(  2,\mathbb{C}%
\right)  .$

\subsubsection{Transpose and Complex Conjugation}

The transpose operation
\begin{align}
\tau &  :\mathcal{B}\left(  \mathcal{H}^{1}\right)  \rightarrow\mathcal{B}%
\left(  \mathcal{H}^{1}\right) \nonumber\\
A^{t}  &  =\tau\left(  A\right)  ,\quad A\in\mathcal{B}\left(  \mathcal{H}%
^{1}\right)
\end{align}
is defined in terms of a complex conjugation map, $K_{\mathfrak{b}},$ that is
an isometric antilinear involutive
\begin{align}
K_{\mathfrak{b}}  &  :\mathcal{H}^{1}\rightarrow\mathcal{H}^{1},\nonumber\\
K_{\mathfrak{b}}\left(  \nu\right)   &  =K_{\mathfrak{b}}\left\vert
\nu\right\rangle =\left\vert \nu^{\ast}\right\rangle ,
\end{align}
that is associated with an equivalence of bases, $\mathfrak{b},$ of
$\mathcal{H}^{1}.$The antilinear property gives
\begin{equation}
K_{\mathfrak{b}}\left(  i\right)  =-i,
\end{equation}
the isometric property
\begin{align}
\left\langle \left.  K_{\mathfrak{b}}\varphi\right\vert K_{\mathfrak{b}}%
\chi\right\rangle  &  =\overline{\left\langle \left.  \varphi\right\vert
\chi\right\rangle }=\left\langle \left.  \chi\right\vert \varphi\right\rangle
\\
&  \Leftrightarrow\left\Vert \varphi\right\Vert =\left\Vert K_{\mathfrak{b}%
}\varphi\right\Vert ,
\end{align}
the involutive property that
\begin{equation}
K_{\mathfrak{b}}^{2}\left\vert \nu\right\rangle =\left\vert \nu\right\rangle ,
\end{equation}
thus
\begin{align}
K_{\mathfrak{b}}\left\vert b_{j}\right\rangle  &  =\left\vert b_{j}%
\right\rangle \nonumber\\
K_{\mathfrak{b}}\left\vert b_{-j}\right\rangle  &  =-\left\vert b_{-j}%
\right\rangle
\end{align}
and
\begin{equation}
\left(  \mathcal{H}^{1}\right)  _{\mathbb{R}}=V_{\mathfrak{b}}\oplus
V_{-\mathfrak{b}}.
\end{equation}
The subspaces $V_{\mathfrak{b}},V_{-\mathfrak{b}}$ are \emph{real }eigenspaces
of $K_{\mathfrak{b}}$.

$\lceil$The adjoint of $K_{\mathfrak{b}}$ is defined as
\begin{equation}
\left\langle \left.  K_{\mathfrak{b}}^{\dagger}\varphi\right\vert
\chi\right\rangle =\left\langle \left.  \varphi\right\vert K_{\mathfrak{b}%
}\chi\right\rangle \quad\forall\left\vert \chi\right\rangle \in\mathcal{H}^{1}%
\end{equation}
$\rfloor$

The matrix representation of an element of $\mathcal{H}^{1}$ and its
$K_{\mathfrak{b}}$-complex conjugate are
\begin{align}
\left\vert v\right\rangle  &  =\left(  \left\vert \lambda_{1}\right\rangle
,\cdots,\left\vert \lambda_{s}\right\rangle ,\left\vert \lambda_{-1}%
\right\rangle ,\cdots,\left\vert \lambda_{-s}\right\rangle \right)
\mathbf{v}\nonumber\\
\left\vert \nu^{\ast}\right\rangle  &  =\left(  \left\vert \lambda
_{1}\right\rangle ,\cdots,\left\vert \lambda_{s}\right\rangle ,\left\vert
\lambda_{-1}\right\rangle ,\cdots,\left\vert \lambda_{-s}\right\rangle
\right)  \mathbf{\bar{v}.}%
\end{align}
If
\begin{equation}
K_{\mathfrak{b}}\left\vert v\right\rangle =\left\vert v\right\rangle
\end{equation}
then $\mathbf{v\in}\mathbb{R}.$

The complex conjugation map, $K_{\mathfrak{b}},$ also defines a non degenerate
symmetric bilinear form on $\mathcal{H}^{1}$ given by
\begin{equation}
\left(  \left.  \varphi\right\vert \chi\right)  _{\mathfrak{b}}=\left(
\left.  \chi\right\vert \varphi\right)  _{\mathfrak{b}}=\left\langle \left.
K_{\mathfrak{b}}\varphi\right\vert \chi\right\rangle =\left\langle \left.
K_{\mathfrak{b}}\chi\right\vert \varphi\right\rangle ;\quad\left\vert
\varphi\right\rangle ,\left\vert \chi\right\rangle \in\mathcal{H}^{1},
\end{equation}
which allows one to define the transpose of an operator as
\begin{equation}
\left\{  \left(  \left.  A^{t}\varphi\right\vert \chi\right)  _{\mathfrak{b}%
}=\left(  \left.  \varphi\right\vert A\chi\right)  _{\mathfrak{b}}\quad
\forall\left\vert \chi\right\rangle \in\mathcal{H}^{1}\right\}  \quad
\forall\left\vert \varphi\right\rangle \in\mathcal{H}^{1},
\end{equation}
hence
\begin{equation}
\left\{  \left(  \left.  A\varphi\right\vert \chi\right)  _{\mathfrak{b}%
}=\left(  \left.  \varphi\right\vert A^{t}\chi\right)  _{\mathfrak{b}}%
\quad\forall\left\vert \chi\right\rangle \in\mathcal{H}^{1}\right\}
\quad\forall\left\vert \varphi\right\rangle \in\mathcal{H}^{1}.
\end{equation}
In particular
\begin{equation}
\left(  \left.  K_{\mathfrak{b}}^{t}\varphi\right\vert \chi\right)
_{\mathfrak{b}}=\left(  \left.  \varphi\right\vert K_{\mathfrak{b}}%
\chi\right)  _{\mathfrak{b}}=\left\langle \left.  K_{\mathfrak{b}}%
\varphi\right\vert K_{\mathfrak{b}}\chi\right\rangle =\overline{\left\langle
\left.  \varphi\right\vert \chi\right\rangle }=\left\langle \left.
\chi\right\vert \varphi\right\rangle ;\quad\left\vert \varphi\right\rangle
,\left\vert \chi\right\rangle \in\mathcal{H}^{1}.
\end{equation}
In terms of the inner product in $\mathcal{H}^{1}$ the transpose definition
can be written as
\begin{equation}
\left\{  \left\langle \left.  K_{\mathfrak{b}}A^{t}K_{\mathfrak{b}%
}K_{\mathfrak{b}}\varphi\right\vert \chi\right\rangle =\left\langle \left.
K_{\mathfrak{b}}\varphi\right\vert A\chi\right\rangle =\left\langle \left.
A^{\dagger}K_{\mathfrak{b}}\varphi\right\vert \chi\right\rangle \quad
\forall\left\vert \chi\right\rangle \in\mathcal{H}^{1}\right\}  \quad
\forall\left\vert \varphi\right\rangle \in\mathcal{H}^{1}%
\end{equation}
which shows that
\begin{equation}
A^{\dagger}=K_{\mathfrak{b}}A^{t}K_{\mathfrak{b}}%
\end{equation}
and
\begin{equation}
K_{\mathfrak{b}}A^{\dagger}K_{\mathfrak{b}}=A^{t}%
\end{equation}
$A$ is $K_{\mathfrak{b}}$-symmetric%
\begin{equation}
\left\{  \left(  \left.  A\varphi\right\vert \chi\right)  _{\mathfrak{b}%
}=\left(  \left.  \varphi\right\vert A\chi\right)  _{\mathfrak{b}}\quad
\forall\left\vert \chi\right\rangle \in\mathcal{H}^{1}\right\}  \quad
\forall\left\vert \varphi\right\rangle \in\mathcal{H}^{1},
\end{equation}
if
\begin{equation}
K_{\mathfrak{b}}A^{\dagger}K_{\mathfrak{b}}=A^{t}=A.
\end{equation}
The change of basis
\begin{align}
S\left(  \left\vert \lambda_{1}\right\rangle ,\cdots,\left\vert \lambda
_{s}\right\rangle ,\left\vert \lambda_{-1}\right\rangle ,\cdots,\left\vert
\lambda_{-s}\right\rangle \right)   &  =\left(  \left\vert \lambda
_{1}\right\rangle ,\cdots,\left\vert \lambda_{s}\right\rangle ,\left\vert
\lambda_{-1}\right\rangle ,\cdots,\left\vert \lambda_{-s}\right\rangle
\right)  \mathbb{S}\nonumber\\
&  =\left(  \left\vert \eta_{1}\right\rangle ,\cdots,\left\vert \eta
_{s}\right\rangle ,\left\vert \eta_{-1}\right\rangle ,\cdots,\left\vert
\eta_{-s}\right\rangle \right)
\end{align}
has real representation matrix $\mathbb{S}$ if $\left\vert \mathbf{\lambda
}\right\rangle $ and $\left\vert \mathbf{\eta}\right\rangle $ belong to the
same $K_{\mathfrak{b}}$-equivalence class i.e. $K_{\lambda}=K_{\eta},$ if they
do not then $\mathbb{S}$ has complex entries.

The complex conjugate of an operator is defined as
\begin{equation}
A^{\ast}=K_{\mathfrak{b}}AK_{\mathfrak{b}},\quad A\in\mathcal{B}\left(
\mathcal{H}^{1}\right)
\end{equation}
and the transpose is then given as
\begin{equation}
A^{t}=\left(  A^{\ast}\right)  ^{\dagger},
\end{equation}
which gives that
\begin{equation}
A^{\dagger}=\left(  A^{t}\right)  ^{\ast}=\left(  A^{\ast}\right)  ^{t}.
\end{equation}
$\lceil$It is possible to show that complex conjugation, transpose and adjoint
operations are associative$\rfloor$ In particular if
\begin{equation}
A=-A^{t}%
\end{equation}
then
\begin{equation}
\left(  iA\right)  ^{\dagger}=-i
\end{equation}


The matrix representation of $A^{\ast}$ on a real $K_{\mathfrak{b}}$-basis is
\begin{equation}
K_{\mathfrak{b}}AK_{\mathfrak{b}}\left(  \left\vert \lambda_{1}\right\rangle
,\cdots,\left\vert \lambda_{s}\right\rangle ,\left\vert \lambda_{-1}%
\right\rangle ,\cdots,\left\vert \lambda_{-s}\right\rangle \right)  =\left(
\left\vert \lambda_{1}\right\rangle ,\cdots,\left\vert \lambda_{s}%
\right\rangle ,\left\vert \lambda_{-1}\right\rangle ,\cdots,\left\vert
\lambda_{-s}\right\rangle \right)  \overline{\mathbb{A}}%
\end{equation}%
\begin{equation}
A^{\ast}=\left\vert \mathbf{\chi}\right\rangle \left\langle \mathbf{\chi
}\left.  A^{\ast}\right\vert \mathbf{\chi}\right\rangle \left\langle
\mathbf{\chi}\right\vert =\left\vert \mathbf{\chi}\right\rangle \left\langle
K_{\mathfrak{b}}\mathbf{\chi}\left.  A\right\vert K_{\mathfrak{b}}%
\mathbf{\chi}\right\rangle \left\langle \mathbf{\chi}\right\vert =\left\vert
\mathbf{\chi}\right\rangle \overline{\mathbb{A}}_{\chi}\left\langle
\mathbf{\chi}\right\vert
\end{equation}
and the representation of $A^{t}$ on such a basis is
\begin{equation}
A^{t}\left(  \left\vert \lambda_{1}\right\rangle ,\cdots,\left\vert
\lambda_{s}\right\rangle ,\left\vert \lambda_{-1}\right\rangle ,\cdots
,\left\vert \lambda_{-s}\right\rangle \right)  =\left(  \left\vert \lambda
_{1}\right\rangle ,\cdots,\left\vert \lambda_{s}\right\rangle ,\left\vert
\lambda_{-1}\right\rangle ,\cdots,\left\vert \lambda_{-s}\right\rangle
\right)  \mathbb{A}^{t}.
\end{equation}


\subsection{Transformations of Symplectic Forms}

The value of the symplectic form $g$ does not depend on its representation i.e.%

\begin{align}
g\left(  \varphi_{1},\varphi_{2}\right)   &  =\left\langle \left.
g\right\vert \varphi_{1}\wedge\varphi_{2}\right\rangle =\left\langle \left.
\sum_{-s\leq k,l\leq s}g_{\chi kl}\left\vert \chi_{k}\chi_{l}\right\rangle
\right\vert \varphi_{1}\wedge\varphi_{2}\right\rangle \nonumber\\
&  =\left\langle \left.  \sum_{1\leq j\leq s}\left\vert c_{j}^{\frac{1}{2}%
}\psi_{j}c_{j}^{\frac{1}{2}}\psi_{-j}\right\rangle \right\vert \varphi
_{1}\wedge\varphi_{2}\right\rangle \nonumber\\
&  =\left\langle \left.  \sum_{1\leq j\leq s}\left\vert \eta_{j}\eta
_{-j}\right\rangle \right\vert \varphi_{1}\wedge\varphi_{2}\right\rangle
;\left\vert \varphi_{1}\right\rangle ,\left\vert \varphi_{2}\right\rangle
\in\mathcal{H}^{1}.
\end{align}


\subsubsection{Discrete Basis}

The canonical form can be obtained by transformations of the one particle,
$\left(  \left\vert \varphi_{1}\right\rangle ,\cdots,\left\vert \varphi
_{s}\right\rangle ,\left\vert \varphi_{-1}\right\rangle ,\cdots,\left\vert
\varphi_{-s}\right\rangle \right)  ,$ of $\mathcal{H}^{1}$ which have the
general effect%
\begin{align}
\left\vert g\right\rangle  &  =\sum_{-s\leq k,l\leq s}g_{\chi kl}\left\vert
\chi_{k}\chi_{l}\right\rangle =\sum_{-s\leq k,l\leq s}g_{\chi kl}\left\vert
T\varphi_{k}T\varphi_{l}\right\rangle =\sum_{-s\leq k,l,k^{\prime},l^{\prime
}\leq s}T_{kk^{\prime}}T_{ll^{\prime}}g_{_{\chi}kl}\left\vert \varphi
_{k}\varphi_{l}\right\rangle \nonumber\\
&  =\sum_{-s\leq k^{\prime},l^{\prime}\leq s}\left\vert \varphi_{k^{\prime}%
}\varphi_{l^{\prime}}\right\rangle \left(  T^{t}\mathbb{G}_{\chi}T\right)
_{k^{\prime}l^{\prime}},
\end{align}
thus
\begin{equation}
\mathbb{T}_{\chi\varphi}^{t}\mathbb{G}_{\varphi}\mathbb{T}_{\chi\varphi
}=\mathbb{G}_{\chi}.
\end{equation}
One also has
\begin{equation}
\mathbb{T}_{\varphi\chi}^{t}\mathbb{G}_{\chi}\mathbb{T}_{\varphi\chi
}=\mathbb{G}_{\varphi}%
\end{equation}
thus
\begin{equation}
\mathbb{T}_{\chi\varphi}^{t}\mathbb{T}_{\varphi\chi}^{t}\mathbb{G}_{\chi
}\mathbb{T}_{\varphi\chi}\mathbb{T}_{\chi\varphi}=\mathbb{G}_{\chi},
\end{equation}
where%
\begin{equation}
\mathbb{G}_{\psi jk}=g\left(  \psi_{j},\psi_{k}\right)  ;j,k=1,\cdots
,s,-1,\cdots,-s.
\end{equation}
The transformation between the canonical and the standard basis is%
\begin{align}
&  \left(
\begin{array}
[c]{cc}%
\mathbb{C}^{-\frac{1}{2}} & \mathbb{O}\\
\mathbb{O} & \mathbb{C}^{-\frac{1}{2}}%
\end{array}
\right)  \left(
\begin{array}
[c]{cc}%
\mathbb{O} & \mathbb{C}\\
-\mathbb{C} & \mathbb{O}%
\end{array}
\right)  \left(
\begin{array}
[c]{cc}%
\mathbb{C}^{-\frac{1}{2}} & \mathbb{O}\\
\mathbb{O} & \mathbb{C}^{-\frac{1}{2}}%
\end{array}
\right)  =\\
&  \left(
\begin{array}
[c]{cc}%
\mathbb{O} & \mathbb{C}^{\frac{1}{2}}\\
-\mathbb{C}^{\frac{1}{2}} & \mathbb{O}%
\end{array}
\right)  \left(
\begin{array}
[c]{cc}%
\mathbb{C}^{-\frac{1}{2}} & \mathbb{O}\\
\mathbb{O} & \mathbb{C}^{-\frac{1}{2}}%
\end{array}
\right)  =\left(
\begin{array}
[c]{cc}%
\mathbb{O} & \mathbb{I}_{s}\\
-\mathbb{I}_{s} & \mathbb{O}%
\end{array}
\right)
\end{align}
thus%
\begin{equation}
\mathbb{T}_{\eta\psi}=\left(
\begin{array}
[c]{cc}%
\mathbb{C}^{-\frac{1}{2}} & \mathbb{O}\\
\mathbb{O} & \mathbb{C}^{-\frac{1}{2}}%
\end{array}
\right)  .
\end{equation}


\subsubsection{Continuous Basis}

The wavefunction associated with $\left\vert g\right\rangle $ is given by
\begin{align}
g\left(  \mathbf{x}_{1},\mathbf{x}_{2}\right)   &  =\left\langle \left.
\mathbf{x}_{1}\mathbf{x}_{2}\right\vert g\right\rangle =\sum_{1\leq j\leq
s}\left\langle \left.  \mathbf{x}_{1}\mathbf{x}_{2}\right\vert \eta_{j}%
\eta_{-j}\right\rangle \\
&  =\sum_{1\leq j\leq s}\left\vert \eta_{j}\left(  \mathbf{x}_{1}\right)
\eta_{-j}\left(  \mathbf{x}_{2}\right)  \right\rangle =\sum_{1\leq j\leq
s}c_{j}\left\vert \psi_{j}\left(  \mathbf{x}_{1}\right)  \psi_{-j}\left(
\mathbf{x}_{2}\right)  \right\rangle ,
\end{align}
thus%
\begin{equation}
c_{j}=\frac{1}{\sqrt{2}}\int\overline{\left\{  \psi_{j}\left(  \mathbf{x}%
_{1}\right)  \psi_{-j}\left(  \mathbf{x}_{2}\right)  -\psi_{j}\left(
\mathbf{x}_{2}\right)  \psi_{-j}\left(  \mathbf{x}_{1}\right)  \right\}
}g\left(  \mathbf{x}_{1},\mathbf{x}_{2}\right)  d\mathbf{x}_{1}d\mathbf{x}_{2}%
\end{equation}
and
\begin{equation}
\int\varphi_{a}\left(  \mathbf{x}_{1}\right)  g\left(  \mathbf{x}%
_{1},\mathbf{x}_{2}\right)  \varphi_{b}\left(  \mathbf{x}_{2}\right)
d\mathbf{x}_{1}d\mathbf{x}_{2}=g\left(  \varphi_{a},\varphi_{b}\right)  ,
\end{equation}
while basis transformations give%
\begin{equation}
g\left(  \mathbf{x}_{1}\mathbf{x}_{2}\right)  =\sum_{-s\leq j\leq s}%
c_{j}\left\vert T\varphi_{j}\left(  \mathbf{x}_{1}\right)  T\varphi
_{-j}\left(  \mathbf{x}_{2}\right)  \right\rangle =\sum_{-s\leq j,k,l\leq
s}T_{jk}T_{-jl}c_{j}\left\vert \varphi_{k}\left(  \mathbf{x}_{1}\right)
\varphi_{l}\left(  \mathbf{x}_{2}\right)  \right\rangle .\nonumber
\end{equation}
The amplitudes, $\left\langle \left.  \mathbf{x}_{1}\right\vert \varphi
_{j}\right\rangle ,$ of the basis elements $\left\vert \varphi_{j}%
\right\rangle $ are the "matrix elements" of the transformation $\Upsilon
:\left(  \left\vert \mathbf{x}_{1}\right\rangle \right)  \rightarrow\left(
\left\vert \varphi_{1}\right\rangle \cdots\left\vert \varphi_{s}\right\rangle
\left\vert \varphi_{-1}\right\rangle \left\vert \varphi_{1}\right\rangle
\left\vert \varphi_{-s}\right\rangle \right)  $ from a discrete basis to a
continuous generalized one i.e.%
\begin{equation}
\left\vert \mathbf{x}_{1}\right\rangle =\sum_{-s\leq j\leq s}\left\vert
\varphi_{j}\right\rangle \left\langle \left.  \varphi_{j}\right\vert
\mathbf{x}_{1}\right\rangle \equiv\sum_{-s\leq j\leq s}\left\vert \varphi
_{j}\right\rangle \overline{\Upsilon_{j}\left(  \mathbf{x}_{1}\right)  }%
\end{equation}
giving
\begin{equation}
\left(  \Upsilon^{\dagger}g\bar{\Upsilon}\right)  \left(  \mathbf{x}%
_{1},\mathbf{x}_{2}\right)  =g\left(  \mathbf{x}_{1},\mathbf{x}_{2}\right)
=\sum_{-s\leq j,k\leq s}\overline{\Upsilon_{j}\left(  \mathbf{x}_{1}\right)
}g_{jk}\overline{\Upsilon_{k}\left(  \mathbf{x}_{2}\right)  }.
\end{equation}
The inverse transformation is given by%
\begin{equation}
\left\vert \varphi_{j}\right\rangle =\int\left\vert \mathbf{x}_{1}%
\right\rangle \left\langle \left.  \mathbf{x}_{1}\right\vert \varphi
_{j}\right\rangle d\mathbf{x}_{1}\equiv\int\left\vert \mathbf{x}%
_{1}\right\rangle \Upsilon_{j}\left(  \mathbf{x}_{1}\right)  d\mathbf{x}_{1}%
\end{equation}
leading to%
\begin{equation}
\left(  \Upsilon^{t}g\Upsilon\right)  _{jk}=\int\Upsilon_{j}\left(
\mathbf{x}_{1}\right)  g\left(  \mathbf{x}_{1}\mathbf{x}_{2}\right)
\Upsilon_{k}\left(  \mathbf{x}_{2}\right)  d\mathbf{x}_{1}d\mathbf{x}_{2}.
\end{equation}


\subsection{Symplectic Group}

The symplectic group $Sp\left(  \mathcal{H}^{1}\right)  \equiv Sp\left(
2s,\mathbb{C}\right)  $ is defined as the set of invertible tranformations $S$
that satisfy
\begin{equation}
g\left(  S\varphi_{1},S\varphi_{2}\right)  =g\left(  \varphi_{1},\varphi
_{2}\right)  \forall\varphi_{1},\varphi_{2}\in\mathcal{H}^{1}%
\end{equation}
if
\begin{equation}
\operatorname*{rank}\left\{  g\right\}  =s.
\end{equation}
If
\begin{equation}
\operatorname*{rank}\left\{  g\right\}  =\rho<s
\end{equation}
then this defines the subgroup $Sp\left(  2\rho,\mathbb{C}\right)  $ of
$Sp\left(  2s,\mathbb{C}\right)  .$

As
\begin{equation}
g\left(  \varphi_{1},\varphi_{2}\right)  =\mathbf{v}_{1\lambda}^{t}%
\mathbb{G}_{\lambda}\mathbf{v}_{2\lambda},
\end{equation}
where
\begin{equation}
\left\vert v\right\rangle =\left(  \left\vert \lambda_{1}\right\rangle
,\cdots,\left\vert \lambda_{s}\right\rangle ,\left\vert \lambda_{-1}%
\right\rangle ,\cdots,\left\vert \lambda_{-s}\right\rangle \right)  \mathbf{v}%
\end{equation}
one has%
\begin{equation}
g\left(  S\varphi_{1},S\varphi_{2}\right)  =\mathbf{v}_{1\lambda}%
^{t}\mathbb{S}_{\lambda}^{t}\mathbb{G}_{\lambda}\mathbb{S}_{\lambda}%
\mathbf{v}_{2\lambda}=\mathbf{v}_{1\lambda}^{t}\mathbb{G}_{\lambda}%
\mathbf{v}_{2\lambda}\quad\forall\mathbf{v}_{1\lambda},\mathbf{v}_{2\lambda
}\in l_{2},
\end{equation}
which shows that%
\begin{equation}
\mathbb{S}_{\lambda}^{t}\mathbb{G}_{\lambda}\mathbb{S}_{\lambda}%
=\mathbb{G}_{\lambda}%
\end{equation}
and
\begin{equation}
\mathbb{S}_{\lambda}^{t}\mathbb{T}_{\lambda\eta}^{t}\mathbb{G}_{\eta
}\mathbb{T}_{\lambda\eta}\mathbb{S}_{\lambda}=\mathbb{G}_{\lambda}%
=\mathbb{T}_{\lambda\eta}^{t}\mathbb{G}_{\eta}\mathbb{T}_{\lambda\eta},
\end{equation}
giving
\begin{align}
&  \mathbb{T}_{\lambda\eta}^{-t}\mathbb{S}_{\lambda}^{t}\mathbb{T}%
_{\eta\lambda}^{t}\mathbb{G}_{\eta}\mathbb{T}_{\eta\lambda}\mathbb{S}%
_{\lambda}\mathbb{T}_{\lambda\eta}^{-1}=\mathbb{G}_{\eta}\nonumber\\
&  \equiv\left[  \mathbb{T}_{\eta\lambda}\mathbb{S}_{\lambda}\mathbb{T}%
_{\lambda\eta}^{-1}\right]  ^{t}\mathbb{G}_{\eta}^{t}\left[  \mathbb{T}%
_{\eta\lambda}\mathbb{S}_{\lambda}\mathbb{T}_{\lambda\eta}^{-1}\right]
=\mathbb{G}_{\eta},
\end{align}
which shows that $\mathbb{S}_{\lambda}$ and $\mathbb{S}_{\eta}$ are conjugate
representations of $Sp\left(  2s,\mathbb{C}\right)  $.

In particular one can choose the biorthogonal dual cannonical basis $\left\{
\left\vert \eta_{j}^{d}\right\rangle ,\left\vert \eta_{-j}^{d}\right\rangle
;1\leq j\leq s\right\}  $ given by
\begin{equation}
\left\vert \eta_{j}^{d}\right\rangle =\left\vert c_{j}^{-\frac{1}{2}}\psi
_{j}\right\rangle ;-s\leq j\leq s
\end{equation}
that produces the form matrix
\begin{equation}
\mathbb{G}_{\eta^{d}}=\left(
\begin{array}
[c]{cc}%
\mathbb{O} & \mathbb{I}\\
-\mathbb{I} & \mathbb{O}%
\end{array}
\right)
\end{equation}
and the group
\begin{equation}
\mathbb{S}_{\eta^{d}}^{t}\mathbb{G}_{\eta^{d}}\mathbb{S}_{\eta^{d}}%
=\mathbb{G}_{\eta^{d}}\equiv\mathbb{S}_{\eta^{d}}^{t}\left(
\begin{array}
[c]{cc}%
\mathbb{O} & \mathbb{I}\\
-\mathbb{I} & \mathbb{O}%
\end{array}
\right)  \mathbb{S}_{\eta^{d}}=\left(
\begin{array}
[c]{cc}%
\mathbb{O} & \mathbb{I}\\
-\mathbb{I} & \mathbb{O}%
\end{array}
\right)  .
\end{equation}
This shows that the invariance group of the general symplectic form is
$Sp\left(  2\rho,\mathbb{C}\right)  .$

\subsection{Invariance Groups}

One can observe from the preceding that the invariance group of $\left\vert
g\right\rangle $ is $Sp\left(  2\rho,\mathbb{C}\right)  ,$ which is a subgroup
of $GL\left(  \mathcal{H}^{1}\right)  $ that has an irreducible representation
on $\mathcal{H}^{2}.$ The hilbert space $\mathcal{H}^{2}$ is not a homogeneous
space for $GL\left(  \mathcal{H}^{1}\right)  ,$ but each orbit does generate
$\mathcal{H}^{2}$ and each vector in $\mathcal{H}^{2}$ is cyclic for this
group and thus all elements of $\mathcal{H}^{2}$ i.e. geminals are generated
by $GL\left(  \mathcal{H}^{1}\right)  $ acting on a reference geminal by the
taking closed linear span of the orbit produced by the action of the coset%
\begin{equation}
\frac{GL\left(  \mathcal{H}^{1}\right)  }{Sp\left(  2\rho,\mathbb{C}\right)
}\equiv\frac{GL\left(  2s,\mathbb{C}\right)  }{Sp\left(  2\rho,\mathbb{C}%
\right)  }%
\end{equation}
on the reference geminal
\begin{equation}
\left\vert g_{\rho}\right\rangle =\sum_{1\leq j\leq s}c_{\rho j}\left\vert
\psi_{\rho j}\psi_{\rho-j}\right\rangle .
\end{equation}
In detail this orbit is given by
\begin{align}
&  \frac{GL\left(  2s,\mathbb{C}\right)  }{Sp\left(  2\rho,\mathbb{C}\right)
}Sp\left(  2\rho,\mathbb{C}\right)  \left\vert g_{\rho}\right\rangle =\left\{
\left.  \left\vert g_{\rho}\left(  \gamma\right)  \right\rangle =\gamma
\left\vert g_{\rho}\right\rangle \right\vert \gamma\in GL\left(
2s,\mathbb{C}\right)  \right\} \nonumber\\
&  =\left\{  \left\vert g_{\rho}\left(  c\right)  \right\rangle =ch\left\vert
g_{\rho}\right\rangle \left\vert c\in\frac{GL\left(  2s,\mathbb{C}\right)
}{Sp\left(  2\rho,\mathbb{C}\right)  },h\in Sp\left(  2\rho,\mathbb{C}\right)
\right.  \right\}  .
\end{align}
The set $\left\{  \left\vert g_{\rho}\left(  c\right)  \right\rangle \right\}
$ is a representation of the homogeneous coset space $\frac{GL\left(
2s,\mathbb{C}\right)  }{Sp\left(  2\rho,\mathbb{C}\right)  }$ as a non linear
manifold contained in $\mathcal{H}^{2}.$

\subsection{Lie Algebra}

The Lie algebra of $Sp\left(  2\rho,\mathbb{C}\right)  $ is generated by
considering one parameter subgroups $\left\{  \left.  S\left(  \tau\right)
\right\vert \tau\in\left(  -\infty,\infty\right)  \right\}  $ of $Sp\left(
2\rho,\mathbb{C}\right)  $
\begin{equation}
S^{t}\left(  \tau\right)  gS\left(  \tau\right)  =g\equiv gS\left(
\tau\right)  =S\left(  \tau\right)  ^{-t}g
\end{equation}
then
\begin{equation}
S\left(  \tau\right)  =e^{s\tau}=I+s\tau+\cdots
\end{equation}
and to first order
\begin{equation}
gs=-s^{t}g=\left(  gs\right)  ^{t}.
\end{equation}
In the standard representation one can find the form of the group matrices $S$
by considering%
\begin{equation}
\left(
\begin{array}
[c]{cc}%
W^{t} & Y^{t}\\
X^{t} & Z^{t}%
\end{array}
\right)  \left(
\begin{array}
[c]{cc}%
\mathbb{O} & \mathbb{I}\\
-\mathbb{I} & \mathbb{O}%
\end{array}
\right)  \left(
\begin{array}
[c]{cc}%
W & X\\
Y & Z
\end{array}
\right)  =\left(
\begin{array}
[c]{cc}%
\mathbb{O} & \mathbb{I}\\
-\mathbb{I} & \mathbb{O}%
\end{array}
\right)
\end{equation}
which leads to%
\begin{equation}
\left(
\begin{array}
[c]{cc}%
-Y^{t} & W^{t}\\
-Z^{t} & X^{t}%
\end{array}
\right)  \left(
\begin{array}
[c]{cc}%
W & X\\
Y & Z
\end{array}
\right)  =\left(
\begin{array}
[c]{cc}%
-Y^{t}W+W^{t}Y & -Y^{t}X+W^{t}Z\\
-Z^{t}W+X^{t}Y & -Z^{t}X+X^{t}Z
\end{array}
\right)  =\left(
\begin{array}
[c]{cc}%
\mathbb{O} & \mathbb{I}\\
-\mathbb{I} & \mathbb{O}%
\end{array}
\right)
\end{equation}
and
\begin{align}
Y^{t}W  &  =\left(  Y^{t}W\right)  ^{t},\\
Z^{t}X  &  =\left(  Z^{t}X\right)  ^{t},\\
-Z^{t}W+X^{t}Y  &  =\mathbb{I}.
\end{align}
The Lie algebra matrices in the standard representation have the following
properties
\begin{equation}
\left(
\begin{array}
[c]{cc}%
\mathbb{O} & \mathbb{I}\\
-\mathbb{I} & \mathbb{O}%
\end{array}
\right)  \left(
\begin{array}
[c]{cc}%
a & b\\
c & d
\end{array}
\right)  =\left(
\begin{array}
[c]{cc}%
c & d\\
-a & -b
\end{array}
\right)  =\left(
\begin{array}
[c]{cc}%
c^{t} & -a^{t}\\
d^{t} & -b^{t}%
\end{array}
\right)
\end{equation}
thus%
\begin{equation}
s=\left(
\begin{array}
[c]{cc}%
a & b\\
c & -a^{t}%
\end{array}
\right)  ;b=b^{t},c=c^{t}.
\end{equation}%
\begin{equation}
gs=-s^{t}g=\left(  gs\right)  ^{t}.
\end{equation}
One can observe that%
\begin{equation}
\left(
\begin{array}
[c]{cc}%
\mathbb{O} & \mathbb{I}\\
-\mathbb{I} & \mathbb{O}%
\end{array}
\right)  \left(
\begin{array}
[c]{cc}%
a & b\\
c & -a^{t}%
\end{array}
\right)  =\left(
\begin{array}
[c]{cc}%
c & -a^{t}\\
-a & -b
\end{array}
\right)  ,
\end{equation}
which is a complex symmetric matrix. The inverse of this transformation is
\begin{equation}
\left(
\begin{array}
[c]{cc}%
\mathbb{O} & \mathbb{-I}\\
\mathbb{I} & \mathbb{O}%
\end{array}
\right)  \left(
\begin{array}
[c]{cc}%
\mathbb{O} & \mathbb{I}\\
-\mathbb{I} & \mathbb{O}%
\end{array}
\right)  =\left(
\begin{array}
[c]{cc}%
\mathbb{I} & \mathbb{O}\\
\mathbb{O} & \mathbb{I}%
\end{array}
\right)
\end{equation}
and%
\begin{equation}
\left(
\begin{array}
[c]{cc}%
-c & a^{t}\\
a & b
\end{array}
\right)  \left(
\begin{array}
[c]{cc}%
\mathbb{O} & \mathbb{I}\\
-\mathbb{I} & \mathbb{O}%
\end{array}
\right)  =\left(
\begin{array}
[c]{cc}%
-a^{t} & -c\\
-b & a
\end{array}
\right)  =-\left(
\begin{array}
[c]{cc}%
a & b\\
c & -a^{t}%
\end{array}
\right)  ^{t}%
\end{equation}
i.e.
\begin{equation}
\left(
\begin{array}
[c]{cc}%
\mathbb{O} & \mathbb{-I}\\
\mathbb{I} & \mathbb{O}%
\end{array}
\right)  \left(
\begin{array}
[c]{cc}%
a & b\\
c & -a^{t}%
\end{array}
\right)  \left(
\begin{array}
[c]{cc}%
\mathbb{O} & \mathbb{I}\\
-\mathbb{I} & \mathbb{O}%
\end{array}
\right)  =-\left(
\begin{array}
[c]{cc}%
a & b\\
c & -a^{t}%
\end{array}
\right)  ^{t}.
\end{equation}
Observing that
\begin{align}
\exp\left\{  TAT^{-1}\right\}   &  =I+TAT^{-1}+\frac{1}{2}TAT^{-1}%
TAT^{-1}+\cdots\\
&  =T\exp\left\{  A\right\}  T^{-1}%
\end{align}
one obtains that
\begin{align}
&  \exp\left\{  \left(
\begin{array}
[c]{cc}%
\mathbb{O} & \mathbb{-I}\\
\mathbb{I} & \mathbb{O}%
\end{array}
\right)  \left(
\begin{array}
[c]{cc}%
a & b\\
c & -a^{t}%
\end{array}
\right)  \left(
\begin{array}
[c]{cc}%
\mathbb{O} & \mathbb{I}\\
-\mathbb{I} & \mathbb{O}%
\end{array}
\right)  \right\} \\
&  =\left(
\begin{array}
[c]{cc}%
\mathbb{O} & \mathbb{-I}\\
\mathbb{I} & \mathbb{O}%
\end{array}
\right)  \exp\left\{  \left(
\begin{array}
[c]{cc}%
a & b\\
c & -a^{t}%
\end{array}
\right)  \right\}  \left(
\begin{array}
[c]{cc}%
\mathbb{O} & \mathbb{I}\\
-\mathbb{I} & \mathbb{O}%
\end{array}
\right)
\end{align}


The canonical representation also has the same structure i.e.
\begin{equation}
\left(
\begin{array}
[c]{cc}%
\mathbb{C}^{\frac{1}{2}} & \mathbb{O}\\
\mathbb{O} & \mathbb{C}^{\frac{1}{2}}%
\end{array}
\right)  \left(
\begin{array}
[c]{cc}%
a & b\\
c & -a^{t}%
\end{array}
\right)  \left(
\begin{array}
[c]{cc}%
\mathbb{C}^{-\frac{1}{2}} & \mathbb{O}\\
\mathbb{O} & \mathbb{C}^{-\frac{1}{2}}%
\end{array}
\right)  =\left(
\begin{array}
[c]{cc}%
a & b\\
c & -a^{t}%
\end{array}
\right)  .
\end{equation}


\subsubsection{Fock Space and Fock Algebra}

The second quantisation map $\Phi:\mathcal{H}^{1}\rightarrow$ $\mathcal{B}%
\left(  \mathcal{H}_{F}\right)  $ is defined as%
\begin{equation}
\Phi\left(  \sum_{1\leq j\leq r}c_{j}\varphi_{j}\right)  =\sum_{1\leq j\leq
r}\bar{c}_{j}\Phi\left(  \varphi_{j}\right)
\end{equation}
and the creators and annihilators associated with the basis $\left\{
\varphi_{j}\right\}  $ as
\begin{align}
a_{j}^{\dagger}  &  =\Phi\left(  \varphi_{j}\right)  ^{\dagger}\nonumber\\
a_{j}  &  =\Phi\left(  \varphi_{j}\right)  .
\end{align}
The space-spin field operators are defined in terms of distributions as
\begin{align}
\Phi\left(  \mathbf{x}\right)   &  =\Phi\left(  \delta_{\mathbf{x}}\right)
=\sum_{1\leq j\leq r}\overline{\varphi_{j}\left(  \mathbf{x}\right)  }%
a_{j}\nonumber\\
\Phi\left(  \mathbf{x}\right)  ^{\dagger}  &  =\Phi\left(  \delta_{\mathbf{x}%
}\right)  ^{\dagger}=\sum_{1\leq j\leq r}\varphi_{j}\left(  \mathbf{x}\right)
a_{j}^{\dagger}%
\end{align}
leading to the geminal expressions%
\begin{equation}
\left\vert g\right\rangle =\sum_{1\leq j\leq s}c_{j}a_{j}^{\dagger}%
a_{-j}^{\dagger}\left\vert \phi\right\rangle =\sum_{1\leq j\leq s}\Phi\left(
\eta_{j}\right)  ^{\dagger}\Phi\left(  \eta_{-j}\right)  ^{\dagger}\left\vert
\phi\right\rangle .
\end{equation}
and%
\begin{equation}
\left\vert g\right\rangle =\int g\left(  \mathbf{x}_{1},\mathbf{x}_{2}\right)
\Phi\left(  \mathbf{x}_{1}\right)  ^{\dagger}\Phi\left(  \mathbf{x}%
_{2}\right)  ^{\dagger}d\mathbf{x}_{1}d\mathbf{x}_{2}\left\vert \phi
\right\rangle
\end{equation}


\subsubsection{Representations}

The Lie algebra generated by
\begin{align}
q_{jk}  &  =c_{j}a_{j}^{\dagger}a_{k}-\sigma\left(  j\right)  \sigma\left(
k\right)  c_{k}a_{-k}^{\dagger}a_{-j};1\leq\left\vert j\right\vert <\left\vert
k\right\vert \leq s,j\neq-k\\
q_{jj}  &  =c_{j}\left(  a_{j}^{\dagger}a_{j}-a_{-j}^{\dagger}a_{-j}\right)
;-s\leq j\leq s\\
q_{j-j}  &  =c_{j}a_{j}^{\dagger}a_{-j};1\leq j\leq s
\end{align}
form a reducible representation $R\left(  sp\left(  2s,\mathbb{C}\right)
\right)  $ of the Lie algebra $sp\left(  2s,\mathbb{C}\right)  $ on the Fock
space, $\mathcal{H}_{F},$ that can be decomposed into the direct sum of
irreps
\begin{equation}
R\left(  sp\left(  2s,\mathbb{C}\right)  \right)  =\bigoplus\limits_{0\leq
N\leq2s}R_{N}\left(  sp\left(  2s,\mathbb{C}\right)  \right)  ,
\end{equation}
where $\left\{  \left.  R_{N}\left(  sp\left(  2s,\mathbb{C}\right)  \right)
\right\vert 0\leq N\leq r\right\}  $ are the irreducible representations
carried by $\mathcal{H}^{N}.$

In general%
\begin{equation}
\Phi\left(  T\varphi_{j}\right)  ^{\dagger}=\Phi\left(  \sum_{1\leq j^{\prime
}\leq r}T_{jj^{\prime}}\varphi_{j^{\prime}}\right)  ^{\dagger}=\sum_{1\leq
j^{\prime}\leq r}T_{jj^{\prime}}\Phi\left(  \varphi_{j^{\prime}}\right)
^{\dagger}=\sum_{1\leq j^{\prime}\leq r}T_{jj^{\prime}}a_{j^{\prime}}%
^{\dagger}%
\end{equation}
leading to
\begin{equation}
\Phi\left(  T\varphi_{j}\right)  ^{\dagger}\Phi\left(  T\varphi_{k}\right)
=\sum_{1\leq j^{\prime},k^{\prime}\leq r}T_{jj^{\prime}}\bar{T}_{kk^{\prime}%
}a_{j^{\prime}}^{\dagger}a_{k^{\prime}}=\sum_{1\leq j^{\prime},k^{\prime}\leq
r}\left(  T\otimes\bar{T}\right)  _{jkj^{\prime}k^{\prime}}a_{j^{\prime}%
}^{\dagger}a_{k^{\prime}}%
\end{equation}
and%
\begin{equation}
Ta_{j}^{\dagger}a_{k}T^{\dagger}=\Phi\left(  T\varphi_{j}\right)  ^{\dagger
}\Phi\left(  T\varphi_{k}\right)  =\sum_{1\leq j^{\prime},k^{\prime}\leq
r}T_{jj^{\prime}}a_{j^{\prime}}^{\dagger}a_{k^{\prime}}T_{k^{\prime}%
k}^{\dagger}.
\end{equation}


\subsection{Lie Algebra Bases}

\subsubsection{Gauss Decomposition Basis}

If the geminal is non degenerate i.e. $c_{j}\neq c_{k}$ the Lie algebra,
$gl\left(  2s,\mathbb{C}\right)  ,$ of $GL\left(  2s,\mathbb{C}\right)  $ is
spanned by the basis elements
\begin{align}
q_{jk}^{\dagger}  &  =\left(  c_{j}a_{k}^{\dagger}a_{j}-\sigma\left(
j\right)  \sigma\left(  k\right)  c_{k}a_{-j}^{\dagger}a_{-k}\right)  ;-s\leq
j<k\leq s,\left\vert j\right\vert \neq\left\vert k\right\vert \nonumber\\
\mathfrak{n}_{j}  &  =c_{j}\left(  a_{j}^{\dagger}a_{j}+a_{-j}^{\dagger}%
a_{-j}\right)  ;1\leq j\leq s\nonumber\\
q_{jk}  &  =\left(  c_{j}a_{j}^{\dagger}a_{k}-\sigma\left(  j\right)
\sigma\left(  k\right)  c_{k}a_{-k}^{\dagger}a_{-j}\right)  ;-s\leq j<k\leq
s;j\neq-k\nonumber\\
q_{j-j}  &  =a_{j}^{\dagger}a_{-j};-s\leq j\leq s,
\end{align}
\emph{which satisfy the indexing symmetry}%
\begin{equation}
q_{jk}=-\sigma\left(  j\right)  \sigma\left(  k\right)  q_{-k-j};1\leq j,k\leq
s.
\end{equation}
In particular
\begin{equation}
q_{jj}=c_{j}\left(  a_{j}^{\dagger}a_{j}-a_{-j}^{\dagger}a_{-j}\right)  ;1\leq
j\leq s\nonumber
\end{equation}
and one has the adjoint relationships%
\begin{align}
q_{j-j}  &  =q_{-jj}^{\dagger},-s\leq j\leq s\nonumber\\
q_{jj}  &  =q_{jj}^{\dagger},\mathfrak{n}_{j}^{\dagger}=\mathfrak{n}_{j},1\leq
j\leq s.
\end{align}
The transformation between the "off diagonal" q operators and the "a"
operators can be succintly displayed for $-s\leq j<k\leq s,\left\vert
k\right\vert \neq\left\vert j\right\vert $ as
\begin{equation}
\left(
\begin{array}
[c]{c}%
q_{jk}\\
q_{j-k}\\
q_{-jk}\\
q_{-j-k}\\
q_{jk}^{\dagger}\\
q_{j-k}^{\dagger}\\
q_{-jk}^{\dagger}\\
q_{-j-k}^{\dagger}%
\end{array}
\right)  =\left(
\begin{array}
[c]{cccccccc}%
c_{j} & 0 & \cdots & \cdots & \cdots & \cdots & 0 & -c_{k}\\
0 & c_{j} & 0 & \cdots & \cdots & 0 & c_{k} & 0\\
\vdots & 0 & c_{j} & 0 & 0 & c_{k} & 0 & \vdots\\
\vdots & \vdots & 0 & c_{j} & -c_{k} & 0 & \vdots & \vdots\\
\vdots & \vdots & 0 & -c_{k} & c_{j} & 0 & \vdots & \vdots\\
\vdots & 0 & c_{k} & 0 & 0 & c_{j} & 0 & \vdots\\
0 & c_{k} & 0 & \cdots & \cdots & 0 & c_{j} & 0\\
-c_{k} & 0 & \cdots & \cdots & \cdots & \cdots & 0 & c_{j}%
\end{array}
\right)  \left(
\begin{array}
[c]{c}%
a_{j}^{\dagger}a_{k}\\
a_{j}^{\dagger}a_{-k}\\
a_{-j}^{\dagger}a_{k}\\
a_{-j}^{\dagger}a_{-k}\\
a_{k}^{\dagger}a_{j}\\
a_{-k}^{\dagger}a_{j}\\
a_{k}^{\dagger}a_{-j}\\
a_{-k}^{\dagger}a_{-j}%
\end{array}
\right)  .
\end{equation}


\subsubsection{Commutation Relationships}

This basis produces a Gauss Decomposition
\begin{equation}
L_{-1}\oplus L_{0}\oplus L_{+1}%
\end{equation}
where
\begin{equation}
L_{+1}=\operatorname*{linspan}\left\{  q_{jk}^{\dagger}=\left(  c_{j}%
a_{k}^{\dagger}a_{j}-\sigma\left(  j\right)  \sigma\left(  k\right)
c_{k}a_{-j}^{\dagger}a_{-k}\right)  ;-s\leq j<k\leq s,\left\vert j\right\vert
\neq\left\vert k\right\vert \right\}
\end{equation}%
\begin{equation}
L_{0}=\operatorname*{linspan}\left\{  q_{j-j}=a_{j}^{\dagger}a_{-j};-s\leq
j\leq s;a_{j}^{\dagger}a_{j}-a_{-j}^{\dagger}a_{-j},1\leq j\leq s\right\}
\end{equation}%
\begin{equation}
L_{-1}=\operatorname*{linspan}\left\{  q_{jk}=\left(  c_{j}a_{j}^{\dagger
}a_{k}-\sigma\left(  j\right)  \sigma\left(  k\right)  c_{k}a_{-k}^{\dagger
}a_{-j}\right)  ;-s\leq j<k\leq s;\left\vert j\right\vert \neq\left\vert
k\right\vert \right\}
\end{equation}
as
\begin{align}
\left[  \mathbf{q},\mathbf{q}\right]   &  \subset L_{-1}\oplus L_{0}\\
\left[  \mathbf{q}^{\dagger},\mathbf{q}^{\dagger}\right]   &  \subset
L_{+1}\oplus L_{0}\\
\left[  \mathbf{q},\mathbf{q}_{0}\right]   &  \subset L_{-1}\\
\left[  \mathbf{q}^{\dagger},\mathbf{q}_{0}\right]   &  \subset L_{+1}\\
\left[  \mathbf{q}_{0},\mathbf{q}_{0}\right]   &  \subset L_{0}%
\end{align}


\subsubsection{Action on Geminal}

The action of these operators on a canonical GSO is%
\begin{equation}
\sum_{1\leq k<l<s}\left\{  z_{kl}q_{kl}^{\dagger}+z_{-kl}q_{-kl}^{\dagger
}+z_{k-l}q_{k-l}^{\dagger}+z_{-k-l}q_{-k-l}^{\dagger}\right\}  \left\vert
\psi_{j}\right\rangle
\end{equation}%
\begin{equation}
z_{kl}=-\sigma\left(  k\right)  \sigma\left(  l\right)  z_{-l-k};-s\leq
k,l\leq s,z_{kl}=0,\left\vert k\right\vert =\left\vert l\right\vert
\end{equation}


The general linear group $GL\left(  \mathcal{H}^{1}\right)  $ and can thus be
parameterized by
\begin{equation}
\gamma\left(  \mathbf{\alpha,z,\beta},\mathbf{w}\right)  =c\left(
\mathbf{\alpha,z}\right)  h\left(  \mathbf{\beta},\mathbf{w}\right)
=c_{1}\left(  \mathbf{\alpha}_{\mathrm{re}},\mathbf{\alpha}_{\mathrm{im}%
}\right)  \tau_{+}\left(  \mathbf{z}\right)  \tau_{-}\left(  \mathbf{w}%
\right)  h_{1}\left(  \mathbf{\beta}\right)
\end{equation}
where
\begin{align}
c\left(  \mathbf{\alpha,z}\right)   &  =c_{1}\left(  \mathbf{\alpha
}_{\mathrm{re}},\mathbf{\alpha}_{\mathrm{im}}\right)  \tau_{+}\left(
\mathbf{z}\right) \nonumber\\
h\left(  \mathbf{\beta},\mathbf{w}\right)   &  =\tau_{-}\left(  \mathbf{w}%
\right)  h_{1}\left(  \mathbf{\beta}\right)
\end{align}
and the parameterization is produced by the coset representatives
\begin{align}
&  \tau_{+}\left(  \mathbf{z}\right)  =\exp\left\{  \sum
\limits_{_{\substack{-s\leq j<k\leq s\\\left\vert j\right\vert \neq\left\vert
k\right\vert }}}z_{jk}q_{jk}^{\dagger}\right\}  ;-s\leq j<k\leq s,z_{jk}%
\in\mathbb{C}\\
&  c_{1i}\left(  \mathbf{\alpha}_{\mathrm{im}}\right)  =\exp\left\{
i\sum\limits_{_{1\leq j\leq s}}\alpha_{\mathrm{im}j}\mathfrak{n}_{j}\right\}
;1\leq j\leq s,\alpha_{\mathrm{im}j}\in\mathbb{R}\\
&  c_{1r}\left(  \mathbf{\alpha}_{\mathrm{re}}\right)  =\exp\left\{
\sum\limits_{_{1\leq j\leq s}}\alpha_{\mathrm{re}j}\mathfrak{n}_{j}\right\}
;1\leq j\leq s,\alpha_{\mathrm{re}j}\in\mathbb{R},
\end{align}
where
\begin{equation}
c_{1}\left(  \mathbf{\alpha}_{\mathrm{re}},\mathbf{\alpha}_{\mathrm{im}%
}\right)  =c_{1i}\left(  \mathbf{\alpha}_{\mathrm{im}}\right)  c_{1r}\left(
\mathbf{\alpha}_{\mathrm{re}}\right)
\end{equation}
and the invariant subgroup
\begin{align}
\tau_{-}\left(  \mathbf{w}\right)   &  =\exp\left\{  \sum
\limits_{_{\substack{-s\leq j<k\leq s\\\left\vert j\right\vert \neq\left\vert
k\right\vert }}}w_{jk}q_{jk}\right\}  ;-s\leq j<k\leq s,w_{jk}\in\mathbb{C}\\
h_{1}\left(  \mathbf{\beta}\right)   &  =\exp\left\{  \sum\limits_{_{1\leq
j\leq s}}\left(  \beta_{jj}\left(  \hat{n}_{j}-\hat{n}_{-j}\right)
+\beta_{j-j}a_{j}^{\dagger}a_{-j}+\beta_{-jj}a_{-j}^{\dagger}a_{j}\right)
\right\}  ,\beta_{jk}\in\mathbb{C}.
\end{align}
The basis $\left\{  \left.  \mathfrak{n}_{j}\right\vert 1\leq j\leq s\right\}
$ for the coset representatives $c_{1}\left(  \mathbf{\alpha}_{\mathrm{re}%
},\mathbf{\alpha}_{\mathrm{im}}\right)  $ can be transformed to
\begin{equation}
\left\{  \left.  \mathcal{N}_{1},\mathfrak{m}_{j}=\mathfrak{n}_{j}%
-\mathfrak{n}_{j+1}\right\vert 1\leq j\leq s-1\right\}  ,
\end{equation}
which allows one to perform the factorization
\begin{equation}
GL\left(  2s,\mathbb{C}\right)  =GL\left(  1,\mathbb{C}\right)  \times
SL\left(  2s,\mathbb{C}\right)
\end{equation}
and set $c\left(  \mathbf{\alpha,z}\right)  $ to
\begin{equation}
c\left(  \mathbf{\tilde{\alpha},z}\right)  =\mathit{s}\left(  \tilde{\alpha
}_{1r}\right)  \mathit{p}\left(  \tilde{\alpha}_{1i}\right)  d\left(
\mathbf{\tilde{\alpha}}\right)  \tau_{+}\left(  \mathbf{z}\right)  ,
\end{equation}
where
\begin{align}
\mathit{s}\left(  \tilde{\alpha}_{1r}\right)  \mathit{p}\left(  \tilde{\alpha
}_{1i}\right)  d\left(  \mathbf{\tilde{\alpha}}\right)   &  =e^{\tilde{\alpha
}_{1r}\mathcal{N}_{1}}e^{i\tilde{\alpha}_{1i}\mathcal{N}_{1}}\exp\left\{
\sum\limits_{_{1\leq j\leq s-1}}\tilde{\alpha}_{j}\mathfrak{m}_{j}\right\} \\
&  \equiv\mathit{s}\left(  \tilde{\alpha}_{1r}\right)  \mathit{p}\left(
\tilde{\alpha}_{1i}\right)  d\left(  \mathbf{\tilde{\alpha}}\right)  ;2\leq
j\leq s,\tilde{\alpha}_{j}\in\mathbb{C},\tilde{\alpha}_{1r},\tilde{\alpha
}_{1i}\in\mathbb{R}%
\end{align}
and $e^{\tilde{\alpha}_{1r}\mathcal{N}_{1}}$ and $e^{i\tilde{\alpha}%
_{1i}\mathcal{N}_{1}}$ are overall scaling and phase factors.
\begin{equation}
\exp\left\{  \sum\limits_{_{1\leq j\leq s-1}}\tilde{\alpha}_{j}\mathfrak{m}%
_{j}\right\}  \exp\left\{  \sum\limits_{_{\substack{-s\leq j<k\leq
s\\\left\vert j\right\vert \neq\left\vert k\right\vert }}}z_{jk}%
q_{jk}^{\dagger}\right\}  \left\vert g_{\rho}\right\rangle =\sum_{1\leq j\leq
s}\left\vert \Upsilon\eta_{\rho j}\Upsilon\eta_{-\rho j}\right\rangle
,\Upsilon\in\frac{SL\left(  2s,\mathbb{C}\right)  }{Sp\left(  2\rho
,\mathbb{C}\right)  }%
\end{equation}


The operators $\left\{  \left.  q_{jk}^{\dagger}\right\vert -s\leq j<k\leq
s,\left\vert j\right\vert \neq\left\vert k\right\vert \right\}  $ and
$\left\{  \left.  q_{jk}\right\vert -s\leq j<k\leq s,\left\vert j\right\vert
\neq\left\vert k\right\vert \right\}  $ satisfy the same commutation
relationships and are representatives of the coset $\frac{Sp\left(
2s,\mathbb{C}\right)  }{SL\left(  2\right)  ^{s}}$ where
\begin{equation}
Sp\left(  2s,\mathbb{C}\right)  =\frac{Sp\left(  2s,\mathbb{C}\right)
}{SL\left(  2\right)  ^{s}}SL\left(  2\right)  ^{s}.
\end{equation}
(It is important to note that in the representation produced by the non
degenerate $q$ operators there is \emph{no} representation of the subgroup
$Sp\left(  2s,\mathbb{C}\right)  ).$ One can observe from the preceeding that
\begin{equation}
SL\left(  2s,\mathbb{C}\right)  =\frac{Sp\left(  2s,\mathbb{C}\right)
}{SL\left(  2\right)  ^{s}}\frac{Sp\left(  2s,\mathbb{C}\right)  }{SL\left(
2\right)  ^{s}}SL\left(  2\right)  ^{s}\cong\frac{Sp\left(  2s,\mathbb{C}%
\right)  }{SL\left(  2\right)  ^{s}}SL\left(  2\right)  ^{s}\frac{Sp\left(
2s,\mathbb{C}\right)  }{SL\left(  2\right)  ^{s}}.
\end{equation}
One thus also has that
\begin{equation}
SL\left(  2s,\mathbb{C}\right)  =Sp\left(  2s,\mathbb{C}\right)  SL\left(
2\right)  ^{s}Sp\left(  2s,\mathbb{C}\right)  ,
\end{equation}
which is the complexification of
\begin{equation}
SU^{\ast}\left(  2s\right)  SU\left(  2\right)  ^{s}SU^{\ast}\left(
2s\right)  .
\end{equation}%
\begin{align}
\lceil\mathfrak{B}_{1}\left(  \mathbb{H}^{1}\right)   &  \cong SU\left(
2\right)  \equiv SU\left(  2,\mathbb{C}\right)  \\
\mathbb{H}^{1} &  =\left\{  \left.  \left(
\begin{array}
[c]{cc}%
x & y\\
-\bar{y} & \overline{x}%
\end{array}
\right)  \right\vert x,y\in\mathbb{C}\right\}  \rfloor
\end{align}


Hence%
\begin{align}
&  \exp\left\{  \sum\limits_{_{\substack{-s\leq j<k\leq s\\\left\vert
j\right\vert \neq\left\vert k\right\vert }}}z_{jk}q_{jk}^{\dagger}\right\}
\left(  \left\vert \psi_{1}\right\rangle \cdots\left\vert \psi_{s}%
\right\rangle ,\left\vert \psi_{-1}\right\rangle \cdots\left\vert \psi
_{-s}\right\rangle \right) \\
&  =\left(  \left\vert \psi_{1}\right\rangle \cdots\left\vert \psi
_{s}\right\rangle ,\left\vert \psi_{-1}\right\rangle \cdots\left\vert
\psi_{-s}\right\rangle \right)  \mathbb{W}\left(  \mathbb{Z}\right)  ,
\end{align}
where $\mathbb{W}\left(  \mathbb{Z}\right)  $ is a matrix representative of
$\frac{Sp\left(  2s,\mathbb{C}\right)  }{SL\left(  2\right)  ^{s}}$ wrt the
basis $\left\{  \left.  \left\vert \psi_{j}\right\rangle \right\vert -s\leq
j\leq s\right\}  $ that depends on $\mathbb{Z}.$

\subsubsection{Cartan Decomposition}

In order to obtain the Cartan decomposition of the Lie algebra one considers
the basis
\begin{align}
u_{jk}  &  =c_{k}a_{j}^{\dagger}a_{k}+\sigma\left(  j\right)  \sigma\left(
k\right)  c_{j}a_{-k}^{\dagger}a_{-j};-s\leq j<k\leq s,\left\vert j\right\vert
\neq\left\vert k\right\vert \\
\mathfrak{m}_{j}  &  =c_{j}\left(  a_{j}^{\dagger}a_{j}+a_{-j}^{\dagger}%
a_{-j}\right)  -c_{j+1}\left(  a_{j+1}^{\dagger}a_{j+1}+a_{-\left(
j+1\right)  }^{\dagger}a_{-\left(  j+1\right)  }\right)  ;1\leq j\leq s-1\\
v_{jk}  &  =c_{k}a_{-k}^{\dagger}a_{-j}-\sigma\left(  j\right)  \sigma\left(
k\right)  c_{j}a_{j}^{\dagger}a_{k};\\
v_{j-j}  &  =a_{j}^{\dagger}a_{-j};-s\leq j\leq s\\
v_{jj}  &  =c_{j}\left(  a_{j}^{\dagger}a_{j}-a_{-j}^{\dagger}a_{-j}\right)
;1\leq j\leq s
\end{align}
The transformation between the "off diagonal" $\left(  q,u\right)  $ operators
and the "a" operators can be succintly displayed for $-s\leq j<k\leq
s,\left\vert k\right\vert \neq\left\vert j\right\vert $ as%
\begin{equation}
\left(
\begin{array}
[c]{c}%
q_{jk}\\
q_{j-k}\\
q_{-jk}\\
q_{-j-k}\\
u_{jk}\\
u_{j-k}\\
u_{-jk}\\
u_{-j-k}%
\end{array}
\right)  =\left(
\begin{array}
[c]{cccccccc}%
c_{j} & 0 & \cdots & \cdots & \cdots & \cdots & 0 & -c_{k}\\
0 & c_{j} & 0 & \cdots & \cdots & 0 & c_{k} & 0\\
\vdots & 0 & c_{j} & 0 & 0 & c_{k} & 0 & \vdots\\
\vdots & \vdots & 0 & c_{j} & -c_{k} & 0 & \vdots & \vdots\\
\vdots & \vdots & 0 & c_{k} & c_{j} & 0 & \vdots & \vdots\\
\vdots & 0 & -c_{k} & 0 & 0 & c_{j} & 0 & \vdots\\
0 & -c_{k} & 0 & \cdots & \cdots & 0 & c_{j} & 0\\
c_{k} & 0 & \cdots & \cdots & \cdots & \cdots & 0 & c_{j}%
\end{array}
\right)  \left(
\begin{array}
[c]{c}%
a_{j}^{\dagger}a_{k}\\
a_{j}^{\dagger}a_{-k}\\
a_{-j}^{\dagger}a_{k}\\
a_{-j}^{\dagger}a_{-k}\\
a_{k}^{\dagger}a_{j}\\
a_{-k}^{\dagger}a_{j}\\
a_{k}^{\dagger}a_{-j}\\
a_{-k}^{\dagger}a_{-j}%
\end{array}
\right)  .
\end{equation}


\subsubsection{Action on Geminal}%

\begin{equation}
u_{kl}\left\vert g\right\rangle =\left(  c_{k}a_{l}^{\dagger}a_{k}+c_{l}%
a_{-k}^{\dagger}a_{-l}\right)  \sum_{1\leq j\leq s}c_{j}\left\vert \psi
_{j}\psi_{-j}\right\rangle =\left(  n_{k}+n_{l}\right)  \left\vert \psi
_{l}\psi_{-k}\right\rangle
\end{equation}%
\begin{equation}
u_{k-l}\left\vert g\right\rangle =\left(  c_{k}a_{-l}^{\dagger}a_{k}%
-c_{l}a_{-k}^{\dagger}a_{l}\right)  \sum_{1\leq j\leq s}c_{j}\left\vert
\psi_{j}\psi_{-j}\right\rangle =\left(  n_{k}+n_{l}\right)  \left\vert
\psi_{-l}\psi_{-k}\right\rangle
\end{equation}%
\begin{equation}
u_{-k-l}\left\vert g\right\rangle =\left(  c_{k}a_{l}^{\dagger}a_{-k}%
-c_{l}a_{k}^{\dagger}a_{-l}\right)  \sum_{1\leq j\leq s}c_{j}\left\vert
\psi_{j}\psi_{-j}\right\rangle =\left(  n_{k}+n_{l}\right)  \left\vert
\psi_{k}\psi_{l}\right\rangle
\end{equation}%
\begin{equation}
u_{-k-l}\left\vert g\right\rangle =\left(  c_{k}a_{-l}^{\dagger}a_{-k}%
+c_{l}a_{k}^{\dagger}a_{l}\right)  \sum_{1\leq j\leq s}c_{j}\left\vert
\psi_{j}\psi_{-j}\right\rangle =\left(  n_{k}+n_{l}\right)  \left\vert
\psi_{k}\psi_{-l}\right\rangle
\end{equation}


\subsubsection{Orbits of $GL\left(  2s,\mathbb{C}\right)  $ generated by
$\left\vert g_{\rho}\right\rangle $}

The orbits of $GL\left(  2s,\mathbb{C}\right)  $ generated by $\left\vert
g_{\rho}\right\rangle $ are given by
\begin{align}
\gamma\left(  \mathbf{\alpha,z,\beta},\mathbf{w}\right)  \left\vert g_{\rho
}\right\rangle  &  =c\left(  \mathbf{\alpha,z}\right)  h\left(  \mathbf{\beta
},\mathbf{w}\right)  \left\vert g_{\rho}\right\rangle =c_{1}\left(
\mathbf{\tilde{\alpha}}\right)  \tau_{+}\left(  \mathbf{z}\right)  \tau
_{-}\left(  \mathbf{w}\right)  h_{1}\left(  \mathbf{\beta}\right)  \left\vert
g_{\rho}\right\rangle \nonumber\\
&  =c_{1}\left(  \mathbf{\tilde{\alpha}}\right)  \tau_{+}\left(
\mathbf{z}\right)  \left\vert g_{\rho}\right\rangle
\end{align}
which gives a representation of the coset space $\frac{GL\left(
1,\mathbb{C}\right)  \times SL\left(  2s,\mathbb{C}\right)  }{sp\left(
2s,\mathbb{C}\right)  }$ as the nonlinear manifold
\begin{equation}
\left\vert g_{\rho}\left(  \mathbf{\tilde{\alpha},z}\right)  \right\rangle
=c_{1}\left(  \mathbf{\tilde{\alpha}}\right)  \tau_{+}\left(  \mathbf{z}%
\right)  \left\vert g_{\rho}\right\rangle =\mathit{s}\left(  \tilde{\alpha
}_{1r}\right)  \mathit{p}\left(  \tilde{\alpha}_{1i}\right)  d\left(
\mathbf{\tilde{\alpha}}\right)  \tau_{+}\left(  \mathbf{z}\right)  \left\vert
g_{\rho}\right\rangle .
\end{equation}
One can restrict the orbit to lie in unit ball of $\mathcal{H}^{2}$ if
$\left\Vert \left\vert g_{\rho}\right\rangle \right\Vert =1$ by setting
\begin{equation}
\tilde{\alpha}_{1r}\left(  \mathbf{\tilde{\alpha},z}\right)  =-\frac{1}{4}%
\ln\left\{  \left\langle g_{\rho}\left(  \mathbf{\tilde{\alpha},z}\right)
\left\vert g_{\rho}\left(  \mathbf{\tilde{\alpha},z}\right)  \right.
\right\rangle \right\}  =-\frac{1}{4}\ln\left\{  \left\langle g_{\rho
}\left\vert d\left(  \mathbf{\tilde{\alpha}}\right)  \tau_{+}\left(
\mathbf{z}\right)  ^{\dagger}\tau_{+}\left(  \mathbf{z}\right)  d\left(
\mathbf{\tilde{\alpha}}\right)  g_{\rho}\right.  \right\rangle \right\}
\end{equation}


\section{Antisymmetrized Geminal Powers (AGP's)}

A reference $2N$-electron AGP\ state in $\mathcal{H}^{2N}$ is defined by
\begin{equation}
\left\vert g_{\rho}^{N}\right\rangle =\sum_{1\leq j_{1}<j_{2}<\cdots<j_{N}\leq
s}c_{\rho j_{1}}\cdots c_{\rho j_{N}}\left\vert \psi_{j_{1}}\psi_{-j_{1}%
}\cdots\psi_{j_{N}\,}\psi_{-j_{N}}\right\rangle \label{canexp}%
\end{equation}
and the wavefunction associated with this state is
\begin{equation}
g_{\rho}^{N}\left(  \mathbf{x}_{1}\cdots\mathbf{x}_{2N}\right)  =\left\langle
\mathbf{x}_{1}\cdots\mathbf{x}_{2N}\left\vert g_{\rho}^{N}\right.
\right\rangle =Pf\left(  g_{\rho}\right)  \left(  \mathbf{x}_{1}%
\cdots\mathbf{x}_{2N}\right)  .
\end{equation}
The CI expansion is given by
\begin{equation}
\operatorname*{Pf}\left(  \Upsilon^{\dagger}g\bar{\Upsilon}\right)  =\sum
\det\left(  \bar{\Upsilon}\right)  _{P}\operatorname*{Pf}\left(  g_{\rho
}\right)  _{P}%
\end{equation}
This reference state can be expressed in operator form as
\begin{align}
\left\vert g_{\rho}^{N}\right\rangle  &  =\sum_{1\leq j_{1}<j_{2}<\cdots
<j_{N}\leq s}c_{\rho j_{1}}\cdots c_{\rho j_{N}}a_{j_{1}}^{\dagger}a_{-j_{1}%
}^{\dagger}\cdots a_{j_{N}}^{\dagger}a_{-j_{N}}^{\dagger}\left\vert
\phi\right\rangle \nonumber\\
&  =\frac{1}{2^{N}N!}\int g_{\rho}\left(  \mathbf{x}_{1},\mathbf{x}%
_{2}\right)  \cdots g_{\rho}\left(  \mathbf{x}_{2N-1},\mathbf{x}_{\not 2
N}\right)  \Phi\left(  \mathbf{x}_{1}\right)  ^{\dagger}\cdots\Phi\left(
\mathbf{x}_{2N}\right)  ^{\dagger}d\mathbf{x}_{1}\cdots d\mathbf{x}%
_{2N}\left\vert \phi\right\rangle \nonumber\\
&  =\frac{1}{2^{N}N!}\int Pf\left(  g_{\rho}\right)  \left(  \mathbf{x}%
_{1}\cdots\mathbf{x}_{2N}\right)  \Phi\left(  \mathbf{x}_{1}\right)
^{\dagger}\cdots\Phi\left(  \mathbf{x}_{2N}\right)  ^{\dagger}d\mathbf{x}%
_{1}\cdots d\mathbf{x}_{2N}\left\vert \phi\right\rangle .
\end{align}
The norm squared of this state is
\begin{align}
\left\langle g_{\rho}^{N}\left\vert g_{\rho}^{N}\right.  \right\rangle  &
=\left(  \frac{1}{2^{N}N!}\right)  ^{2}\int\overline{Pf\left(  g_{\rho
}\right)  \left(  \mathbf{x}_{1}\cdots\mathbf{x}_{2N}\right)  }Pf\left(
g_{\rho}\right)  \left(  \mathbf{x}_{1}\cdots\mathbf{x}_{2N}\right)
d\mathbf{x}_{1}\cdots d\mathbf{x}_{2N}\nonumber\\
&  =\left(  \frac{1}{2^{N}N!}\right)  ^{2}\int\left\vert Pf\left(  g_{\rho
}\right)  \left(  \mathbf{x}_{1}\cdots\mathbf{x}_{2N}\right)  \right\vert
^{2}d\mathbf{x}_{1}\cdots d\mathbf{x}_{2N}\nonumber\\
&  =\left(  \frac{1}{2^{N}N!}\right)  ^{2}\int\left\vert \det\left(  g_{\rho
}\right)  \left(  \mathbf{x}_{1}\cdots\mathbf{x}_{2N}\right)  \right\vert
d\mathbf{x}_{1}\cdots d\mathbf{x}_{2N},
\end{align}
where%
\begin{equation}
\det\left(  g_{\rho}\right)  \left(  \mathbf{x}_{1}\cdots\mathbf{x}%
_{2N}\right)  =\det\left(
\begin{array}
[c]{ccc}%
g_{\rho}\left(  \mathbf{x}_{1},\mathbf{x}_{1}\right)  & \cdots & g_{\rho
}\left(  \mathbf{x}_{1},\mathbf{x}_{2N}\right) \\
\vdots & \ddots & \vdots\\
g_{\rho}\left(  \mathbf{x}_{2N},\mathbf{x}_{1}\right)  & \cdots & g_{\rho
}\left(  \mathbf{x}_{2N},\mathbf{x}_{2N}\right)
\end{array}
\right)  .
\end{equation}
This should be compared to
\begin{equation}
\Phi\left(  \mathbf{x}_{1}\cdots\mathbf{x}_{2N}\right)  =\det\left(
\begin{array}
[c]{ccc}%
\varphi_{1}\left(  \mathbf{x}_{1}\right)  & \cdots & \varphi_{2N}\left(
\mathbf{x}_{1}\right) \\
\vdots & \ddots & \vdots\\
\varphi_{1}\left(  \mathbf{x}_{2N}\right)  & \cdots & \varphi_{2N}\left(
\mathbf{x}_{2N}\right)
\end{array}
\right)
\end{equation}
and%
\begin{align}
&  \left\langle \Phi\left\vert \Phi\right.  \right\rangle =\\
&  \int\overline{\det\left(
\begin{array}
[c]{ccc}%
\varphi_{1}\left(  \mathbf{x}_{1}\right)  & \cdots & \varphi_{2N}\left(
\mathbf{x}_{1}\right) \\
\vdots & \ddots & \vdots\\
\varphi_{1}\left(  \mathbf{x}_{2N}\right)  & \cdots & \varphi_{2N}\left(
\mathbf{x}_{2N}\right)
\end{array}
\right)  }\det\left(
\begin{array}
[c]{ccc}%
\varphi_{1}\left(  \mathbf{x}_{1}\right)  & \cdots & \varphi_{2N}\left(
\mathbf{x}_{1}\right) \\
\vdots & \ddots & \vdots\\
\varphi_{1}\left(  \mathbf{x}_{2N}\right)  & \cdots & \varphi_{2N}\left(
\mathbf{x}_{2N}\right)
\end{array}
\right)  d\mathbf{x}_{1}\cdots d\mathbf{x}_{2N}\\
&  =\det\left(
\begin{array}
[c]{ccc}%
\left\langle \varphi_{1}\left\vert \varphi_{1}\right.  \right\rangle  & \cdots
& \left\langle \varphi_{1}\left\vert \varphi_{2N}\right.  \right\rangle \\
\vdots & \ddots & \vdots\\
\left\langle \varphi_{2N}\left\vert \varphi_{1}\right.  \right\rangle  &
\cdots & \left\langle \varphi_{2N}\left\vert \varphi_{2N}\right.
\right\rangle
\end{array}
\right)  .
\end{align}


The invariance group of $\left\vert g_{\rho}^{N}\right\rangle $ is again
$sp\left(  2s,\mathbb{C}\right)  $ and the orbit of $\left\vert g_{\rho}%
^{N}\right\rangle $ under $GL\left(  2s,\mathbb{C}\right)  $ produces a
representation of the coset $\frac{GL\left(  1,\mathbb{C}\right)  \times
SL\left(  2s,\mathbb{C}\right)  }{sp\left(  2s,\mathbb{C}\right)  }$ in
$\mathcal{H}^{2N}$ by the manifold
\begin{equation}
\left\vert g_{\rho}\left(  \mathbf{\tilde{\alpha},z}\right)  ^{N}\right\rangle
=c_{1}\left(  \mathbf{\tilde{\alpha}}\right)  \tau_{+}\left(  \mathbf{z}%
\right)  \left\vert g_{\rho}^{N}\right\rangle
\end{equation}
which is parameterized by $\left\{  \mathbf{\tilde{\alpha},z}\right\}  .$ One
can consider the manifold of normalized AGP's by restricting%
\begin{align}
\tilde{\alpha}_{1r}\left(  \mathbf{\tilde{\alpha},z}\right)   &  =-\frac
{1}{4N}\ln\left\{  \left\langle g_{\rho}\left(  \mathbf{\tilde{\alpha}%
,z}\right)  ^{N}\left\vert g_{\rho}\left(  \mathbf{\tilde{\alpha},z}\right)
^{N}\right.  \right\rangle \right\} \nonumber\\
&  =-\frac{1}{4N}\ln\left\{  \left\langle g_{\rho}^{N}\left\vert c\left(
\mathbf{\tilde{\alpha}}\right)  \tau_{+}\left(  \mathbf{z}\right)  ^{\dagger
}\tau_{+}\left(  \mathbf{z}\right)  c\left(  \mathbf{\tilde{\alpha}}\right)
g_{\rho}^{N}\right.  \right\rangle \right\} \nonumber\\
&  =-\frac{1}{4N}\ln\left\{  S_{N}\left(  \mathbf{\tilde{\alpha}}\right)
\right\}  .
\end{align}


The manifold generated by a reference AGP state of a given rank is an orbit of
the group $GL\left(  2s,\mathbb{C}\right)  $ and the group is transitive on
that orbit (we can take $\rho$ to label rank). The linear span of any such
orbit is the whole space $\mathcal{H}^{2N}$ thus any reference AGP vector in
$\mathcal{H}^{2N}$ is cyclic for $GL\left(  2s,\mathbb{C}\right)  $.

\section{Action of the group on GSO's}%

\begin{equation}
\gamma\left\vert g\right\rangle =\gamma\sum_{1\leq j\leq s}c_{j}\left\vert
\eta_{-j}\eta_{j}\right\rangle =\sum_{1\leq j\leq s}c_{j}\left\{  \left\vert
\left(  \gamma\eta_{-j}\right)  \eta_{j}\right\rangle +\left\vert \eta
_{-j}\left(  \gamma\eta_{j}\right)  \right\rangle \right\}
\end{equation}%
\begin{equation}
\gamma\left\vert \eta_{j}\right\rangle =\gamma\left(  \mathbf{z,\eta
,w}\right)  \left\vert \eta_{j}\right\rangle =c\left(  \mathbf{\eta
}_{\mathrm{re}}\right)  c\left(  \mathbf{\eta}_{\mathrm{im}}\right)  \tau
_{+}\left(  \mathbf{z}\right)  \left\vert \eta_{j}\right\rangle
\end{equation}%
\begin{align}
\tau_{+}\left(  \mathbf{z}\right)  \left\vert \eta_{j}\right\rangle  &
=\exp\left\{  \sum\limits_{_{\substack{-s\leq k<l\leq s\\\left\vert
k\right\vert \neq\left\vert l\right\vert }}}z_{kl}q_{jl}^{\dagger}\right\}
\left\vert \eta_{j}\right\rangle \\
&  =\exp\left\{  \sum\limits_{_{\substack{-s\leq k<l\leq s\\\left\vert
k\right\vert \neq\left\vert l\right\vert }}}a_{p}^{\dagger}a_{q}T_{pqkl}%
z_{kl}\right\}  \left\vert \eta_{j}\right\rangle
\end{align}


\section{Example of Indexing}

The operators are indexed as%

\begin{equation}
q_{jk}=\left(  c_{j}a_{j}^{\dagger}a_{k}-\sigma\left(  j\right)  \sigma\left(
k\right)  c_{k}a_{-k}^{\dagger}a_{-j}\right)  ;-s\leq j<k\leq s;-s\leq
\left\vert j\right\vert \neq\left\vert k\right\vert \leq s
\end{equation}
with the symmetry%
\begin{align}
q_{jk}  &  =\left(  c_{j}a_{j}^{\dagger}a_{k}-\sigma\left(  j\right)
\sigma\left(  k\right)  c_{k}a_{-k}^{\dagger}a_{-j}\right)  =\\
-\sigma\left(  jk\right)  q_{-k-j}  &  =-\sigma\left(  jk\right)  \left(
c_{k}a_{-k}^{\dagger}a_{-j}-\sigma\left(  j\right)  \sigma\left(  k\right)
c_{j}a_{j}^{\dagger}a_{k}\right)
\end{align}
and
\begin{equation}
q_{jk}^{\dagger}=\left(  c_{j}a_{k}^{\dagger}a_{j}-\sigma\left(  j\right)
\sigma\left(  k\right)  c_{k}a_{-j}^{\dagger}a_{-k}\right)  ;-s\leq j<k\leq
s;-s\leq\left\vert j\right\vert \neq\left\vert k\right\vert \leq s.
\end{equation}
The diagonal blocks are given by
\begin{align}
q_{j-j}  &  =a_{j}^{\dagger}a_{-j}=q_{-jj}^{\dagger}\\
q_{jj}  &  =-q_{-j-j}=c_{j}\left(  a_{j}^{\dagger}a_{j}-a_{-j}^{\dagger}%
a_{-j}\right)  =q_{jj}^{\dagger}=-q_{-j-j}^{\dagger}.
\end{align}
for $1\leq j\leq s.$

\subsection{$sp(4,\mathbb{C)}$}

The following example demonstrates the indexing for $sp(4,\mathbb{C)}.$ The
basic pattern is%
\begin{equation}
\left(  \left\vert \psi_{1}\right\rangle \left\vert \psi_{2}\right\rangle
\left\vert \psi_{-1}\right\rangle \left\vert \psi_{-2}\right\rangle \right)
\left(
\begin{array}
[c]{cccc}%
a & b & e & f\\
c & d & f & g\\
h & j & -a & -c\\
j & k & -b & -d
\end{array}
\right)  \left(
\begin{array}
[c]{c}%
\left\langle \psi_{1}\right\vert \\
\left\langle \psi_{2}\right\vert \\
\left\langle \psi_{-1}\right\vert \\
\left\langle \psi_{-2}\right\vert
\end{array}
\right)  .
\end{equation}
The off diagonal "block" elements are
\begin{equation}
q_{-22},q_{2-2},q_{-11},q_{1-1},
\end{equation}
where
\begin{align}
q_{-22}  &  =c_{2}a_{2}^{\dagger}a_{-2},q_{2-2}=c_{2}a_{-2}^{\dagger}%
a_{2}\nonumber\\
q_{-11}  &  =c_{1}a_{1}^{\dagger}a_{-1},q_{1-1}=c_{1}a_{-1}^{\dagger}a_{1}%
\end{align}
and the diagonal ones are
\begin{align}
q_{-2-2}  &  =c_{2}\left(  a_{-2}^{\dagger}a_{-2}-a_{2}^{\dagger}a_{2}\right)
=-q_{22}\nonumber\\
q_{-1-1}  &  =c_{1}\left(  a_{-1}^{\dagger}a_{-1}-a_{1}^{\dagger}a_{1}\right)
=-q_{11}.
\end{align}
The inter block basis elements are
\begin{align}
q_{12}  &  =c_{1}a_{1}^{\dagger}a_{2}-c_{2}a_{-2}^{\dagger}a_{-1}%
=-q_{-2-1}\nonumber\\
q_{-12}  &  =c_{1}a_{-1}^{\dagger}a_{2}+c_{2}a_{-2}^{\dagger}a_{1}%
=q_{-21}\nonumber\\
q_{1-2}  &  =c_{1}a_{1}^{\dagger}a_{-2}+c_{2}a_{2}^{\dagger}a_{-1}%
=q_{2-1}\nonumber\\
q_{-1-2}  &  =c_{1}a_{-1}^{\dagger}a_{-2}-c_{2}a_{2}^{\dagger}a_{1}=-q_{21}.
\end{align}


The commutation relationships based on
\begin{align}
\left[  q_{ij},q_{kl}\right]   &  =\sigma\left(  jk\right)  \left\{
c_{j}c_{k}q_{il}\delta_{kj}+c_{i}c_{l}q_{-j-k}\delta_{il}+c_{j}c_{l}%
q_{i-k}\delta_{-jl}+c_{i}c_{k}q_{-jl}\delta_{k-i}\right\}  ;\\
-s  &  \leq i,j\leq s;-s\leq k,l\leq s.
\end{align}
Between inter block elements they are given by%
\begin{align}
&  \left[  q_{12},q_{-12}\right]  =\left[  q_{-21},q_{12}\right]  =\left[
\left(  c_{1}a_{1}^{\dagger}a_{2}-c_{2}a_{-2}^{\dagger}a_{-1}\right)  ,\left(
c_{1}a_{-1}^{\dagger}a_{2}+c_{2}a_{-2}^{\dagger}a_{1}\right)  \right]  =\\
&  =-2c_{1}c_{2}a_{-2}^{\dagger}a_{2}=2c_{1}c_{2}q_{-22}\\
&  \left[  q_{12},q_{1-2}\right]  =\left[  \left(  c_{1}a_{1}^{\dagger}%
a_{2}-c_{2}a_{-2}^{\dagger}a_{-1}\right)  ,\left(  c_{1}a_{1}^{\dagger}%
a_{-2}+c_{2}a_{2}^{\dagger}a_{-1}\right)  \right]  =\\
&  =2c_{1}c_{2}a_{1}^{\dagger}a_{-1}=2c_{1}c_{2}q_{1-1}\\
&  \left[  q_{12},q_{-1-2}\right]  =-\left[  q_{12},q_{21}\right]  =\left[
\left(  c_{1}a_{1}^{\dagger}a_{2}-c_{2}a_{-2}^{\dagger}a_{-1}\right)  ,\left(
c_{1}a_{-1}^{\dagger}a_{-2}-c_{2}a_{2}^{\dagger}a_{1}\right)  \right]  =\\
&  =c_{1}c_{2}\left\{  a_{2}^{\dagger}a_{2}-a_{-2}^{\dagger}a_{-2}%
+a_{-1}^{\dagger}a_{-1}-a_{1}^{\dagger}a_{1}\right\}  =c_{1}c_{2}\left\{
q_{22}-q_{11}\right\}
\end{align}%
\begin{align}
&  \left[  q_{-21},q_{1-2}\right]  =\left[  \left(  c_{2}a_{-2}^{\dagger}%
a_{1}+c_{1}a_{-1}^{\dagger}a_{2}\right)  ,\left(  c_{1}a_{1}^{\dagger}%
a_{-2}+c_{2}a_{2}^{\dagger}a_{-1}\right)  \right]  =\nonumber\\
&  c_{1}c_{2}a_{-2}^{\dagger}a_{-2}-c_{1}c_{2}a_{-1}^{\dagger}a_{-1}%
+c_{1}c_{2}a_{-1}^{\dagger}a_{-1}-c_{1}c_{2}a_{-2}^{\dagger}a_{-2}%
=0\nonumber\\
&  \left[  q_{-21},q_{-1-2}\right]  =\left[  \left(  c_{2}a_{-2}^{\dagger
}a_{1}+c_{1}a_{-1}^{\dagger}a_{2}\right)  ,\left(  c_{1}a_{-1}^{\dagger}%
a_{-2}-c_{2}a_{2}^{\dagger}a_{1}\right)  \right]  =\nonumber\\
&  c_{2}c_{1}a_{-1}^{\dagger}a_{1}-c_{2}c_{1}a_{-1}^{\dagger}a_{1}=0
\end{align}%
\begin{align}
&  \left[  q_{1-2},q_{-1-2}\right]  =\left[  q_{21},q_{1-2}\right]  =\left[
\left(  c_{1}a_{1}^{\dagger}a_{-2}+c_{2}a_{2}^{\dagger}a_{-1}\right)  ,\left(
c_{1}a_{-1}^{\dagger}a_{-2}-c_{2}a_{2}^{\dagger}a_{1}\right)  \right]
=\nonumber\\
&  \left.  =\right.  2c_{1}c_{2}a_{2}^{\dagger}a_{-2}=2c_{1}c_{2}q_{2-2}%
\end{align}%
\begin{align}
\left[  q_{12},a_{2}^{\dagger}a_{-2}\right]   &  =\left[  q_{12}%
,q_{2-2}\right]  =\left[  \left(  c_{1}a_{1}^{\dagger}a_{2}-c_{2}%
a_{-2}^{\dagger}a_{-1}\right)  ,a_{2}^{\dagger}a_{-2}\right]  =\nonumber\\
c_{1}a_{1}^{\dagger}a_{-2}+c_{2}a_{2}^{\dagger}a_{-1}  &  =q_{1-2}\nonumber\\
\left[  q_{12},a_{1}^{\dagger}a_{-1}\right]   &  =\left[  q_{12}%
,q_{1-1}\right]  =\left[  \left(  c_{1}a_{1}^{\dagger}a_{2}-c_{2}%
a_{-2}^{\dagger}a_{-1}\right)  ,a_{1}^{\dagger}a_{-1}\right]  =0\nonumber\\
\left[  q_{12},\left(  a_{1}^{\dagger}a_{1}-a_{-1}^{\dagger}a_{-1}\right)
\right]   &  =\left[  \left(  c_{1}a_{1}^{\dagger}a_{2}-c_{2}a_{-2}^{\dagger
}a_{-1}\right)  ,\left(  a_{1}^{\dagger}a_{1}-a_{-1}^{\dagger}a_{-1}\right)
\right]  =0\nonumber\\
\left[  q_{12},\left(  a_{2}^{\dagger}a_{2}-a_{-2}^{\dagger}a_{-2}\right)
\right]   &  =\left[  q_{12},q_{22}\right]  =\left[  \left(  c_{1}%
a_{1}^{\dagger}a_{2}-c_{2}a_{-2}^{\dagger}a_{-1}\right)  ,\left(
a_{2}^{\dagger}a_{2}-a_{-2}^{\dagger}a_{-2}\right)  \right] \nonumber\\
&  =c_{1}a_{1}^{\dagger}a_{2}-c_{2}a_{-2}^{\dagger}a_{-1}=q_{12}%
\end{align}%
\begin{align}
\left[  a_{1}^{\dagger}a_{-1},\left(  a_{1}^{\dagger}a_{1}-a_{-1}^{\dagger
}a_{-1}\right)  \right]   &  =-2a_{1}^{\dagger}a_{-1}\\
\left[  a_{2}^{\dagger}a_{-2},\left(  a_{2}^{\dagger}a_{2}-a_{-2}^{\dagger
}a_{-2}\right)  \right]   &  =-2a_{2}^{\dagger}a_{-2}%
\end{align}


\subsection{$sp(6,\mathbb{C)}$}

The following example demonstrates the indexing for $sp(6,\mathbb{C)}$%
\begin{equation}%
\begin{array}
[c]{ccccccc}
& 1 & 2 & 3 & -1 & -2 & -3\\
1 & a & b & c & m & n & o\\
2 & d & e & f & n & p & q\\
3 & j & k & l & o & q & r\\
-1 & \alpha & \beta & \gamma & -a & -d & -j\\
-2 & \beta & \delta & \varepsilon & -b & -e & -k\\
-3 & \gamma & \varepsilon & \zeta & -c & -f & -l
\end{array}
\end{equation}
The "block diagonal" operators are
\begin{equation}
q_{-3-3},q_{-2-2},q_{-1-1};
\end{equation}
where
\begin{align}
q_{-3-3}  &  =c_{3}\left(  a_{-3}^{\dagger}a_{-3}-a_{3}^{\dagger}a_{3}\right)
\nonumber\\
q_{-2-2}  &  =c_{2}\left(  a_{-2}^{\dagger}a_{-2}-a_{2}^{\dagger}a_{2}\right)
\nonumber\\
q_{-1-1}  &  =c_{1}\left(  a_{-1}^{\dagger}a_{-1}-a_{1}^{\dagger}a_{1}\right)
\end{align}
and the off diagonal operators are
\begin{equation}
q_{-33},q_{3-3},q_{-22},q_{2-2},q_{-11},q_{1-1},
\end{equation}
where%
\begin{align}
q_{-33}  &  =a_{3}^{\dagger}a_{-3};q_{3-3}=a_{-3}^{\dagger}a_{3}\nonumber\\
q_{-22}  &  =a_{2}^{\dagger}a_{-2};q_{2-2}=a_{-2}^{\dagger}a_{2}\nonumber\\
q_{-11}  &  =a_{1}^{\dagger}a_{-1};q_{1-1}=a_{-1}^{\dagger}a_{1}.
\end{align}%
\begin{align}
\left[  q_{ij},q_{kl}\right]   &  =\sigma\left(  jk\right)  \left\{
c_{i}q_{il}\delta_{kj}+c_{i}c_{i}q_{-j-k}\delta_{il}+c_{i}c_{i}q_{i-k}%
\delta_{-jl}+c_{i}c_{i}q_{-jl}\delta_{k-i}\right\}  ;\\
-s  &  \leq i<j\leq s;-s\leq k<l\leq s
\end{align}
The commutation relationships between the block operators are
\begin{align}
\left[  q_{-3-3},q_{-33}\right]   &  =-\left[  c_{3}\left(  a_{3}^{\dagger
}a_{3}-a_{-3}^{\dagger}a_{-3}\right)  ,a_{3}^{\dagger}a_{-3}\right]
=-c_{3}\left(  a_{3}^{\dagger}a_{3}a_{3}^{\dagger}a_{-3}+a_{3}^{\dagger}%
a_{-3}a_{3}^{\dagger}a_{3}\right) \nonumber\\
&  =-c_{3}a_{3}^{\dagger}a_{-3}=c_{3}q_{-33},
\end{align}%
\begin{align}
&  \left[  q_{2-3},q_{-3-3}\right]  =\nonumber\\
&  -\left[  c_{3}\left(  a_{3}^{\dagger}a_{3}-a_{-3}^{\dagger}a_{-3}\right)
,\left(  c_{3}a_{-3}^{\dagger}a_{2}-\sigma\left(  -3\right)  \sigma\left(
2\right)  c_{2}a_{-2}^{\dagger}a_{3}\right)  \right] \nonumber\\
&  =c_{3}^{2}a_{-2}^{\dagger}a_{3}-c_{2}c_{3}a_{-3}^{\dagger}a_{2}%
=c_{3}q_{2-3}.
\end{align}
The off block operators are
\begin{equation}
q_{12},q_{1-2},q_{13},q_{1-3},q_{-12},q_{-1-2},q_{-13},q_{-13},q_{23}%
,q_{2-3},q_{-23},q_{-2-3}%
\end{equation}
are indexed by%
\begin{equation}
\left\{  q_{jk}\left\vert 1\leq\left\vert j\right\vert <\left\vert
k\right\vert \leq s\right.  -s\leq j,k\leq s\right\}  ,
\end{equation}
where
\begin{align}
q_{12}  &  =c_{1}a_{1}^{\dagger}a_{2}-c_{2}a_{-2}^{\dagger}a_{-1}\nonumber\\
q_{1-2}  &  =c_{1}a_{1}^{\dagger}a_{-2}+c_{2}a_{2}^{\dagger}a_{-1}\nonumber\\
q_{-12}  &  =c_{1}a_{-1}^{\dagger}a_{2}+c_{2}a_{-2}^{\dagger}a_{1}\nonumber\\
q_{-1-2}  &  =c_{1}a_{-1}^{\dagger}a_{-2}-c_{2}a_{2}^{\dagger}a_{1}\nonumber\\
q_{13}  &  =c_{1}a_{1}^{\dagger}a_{3}-c_{3}a_{-3}^{\dagger}a_{-1}\nonumber\\
q_{1-3}  &  =c_{1}a_{1}^{\dagger}a_{-3}+c_{3}a_{3}^{\dagger}a_{-1}\nonumber\\
q_{-13}  &  =c_{1}a_{-1}^{\dagger}a_{3}+c_{3}a_{-3}^{\dagger}a_{1}\nonumber\\
q_{-1-3}  &  =c_{1}a_{-1}^{\dagger}a_{-3}-c_{3}a_{3}^{\dagger}a_{1}\nonumber\\
q_{23}  &  =c_{2}a_{2}^{\dagger}a_{3}-c_{3}a_{3}^{\dagger}a_{2}\nonumber\\
q_{2-3}  &  =c_{2}a_{2}^{\dagger}a_{-3}+c_{3}a_{3}^{\dagger}a_{-2}\nonumber\\
q_{-23}  &  =c_{2}a_{-2}^{\dagger}a_{3}+c_{3}a_{-3}^{\dagger}a_{2}\nonumber\\
q_{-2-3}  &  =c_{2}a_{-2}^{\dagger}a_{-3}-c_{3}a_{3}^{\dagger}a_{2}%
\end{align}
The commutation relationships are based on%
\begin{equation}
\left[  a_{r}^{\dagger}a_{s},a_{p}^{\dagger}a_{q}\right]  =a_{r}^{\dagger
}a_{q}\delta_{sp}-a_{p}^{\dagger}a_{s}\delta_{rq}%
\end{equation}
and are given by%
\begin{align}
&  \left[  q_{12},q_{13}\right]  =\left[  c_{1}a_{1}^{\dagger}a_{2}%
-c_{2}a_{-2}^{\dagger}a_{-1},c_{1}a_{1}^{\dagger}a_{3}-c_{3}a_{-3}^{\dagger
}a_{-1}\right]  =0\nonumber\\
&  \left[  q_{12},q_{1-3}\right]  =\left[  c_{1}a_{1}^{\dagger}a_{2}%
-c_{2}a_{-2}^{\dagger}a_{-1},c_{1}a_{1}^{\dagger}a_{-3}+c_{3}a_{3}^{\dagger
}a_{-1}\right]  =0\nonumber\\
&  \left[  q_{12},q_{-13}\right]  =\left[  c_{1}a_{1}^{\dagger}a_{2}%
-c_{2}a_{-2}^{\dagger}a_{-1},c_{1}a_{-1}^{\dagger}a_{3}-c_{3}a_{-3}^{\dagger
}a_{1}\right]  =\nonumber\\
&  c_{1}c_{3}a_{-3}^{\dagger}a_{2}-c_{2}c_{1}a_{-2}^{\dagger}a_{3}%
=c_{1}\left(  c_{3}a_{-3}^{\dagger}a_{2}-c_{2}a_{-2}^{\dagger}a_{3}\right)
=c_{1}q_{-32}\nonumber\\
&  \left[  q_{12},q_{-1-3}\right]  =\left[  c_{1}a_{1}^{\dagger}a_{2}%
-c_{2}a_{-2}^{\dagger}a_{-1},c_{1}a_{-1}^{\dagger}a_{-3}-c_{3}a_{3}^{\dagger
}a_{1}\right]  =\nonumber\\
&  c_{1}c_{3}a_{3}^{\dagger}a_{2}-c_{1}c_{2}a_{-2}^{\dagger}a_{-3}%
=c_{1}\left(  c_{3}a_{3}^{\dagger}a_{2}-c_{2}a_{-2}^{\dagger}a_{-3}\right)
=c_{1}q_{32}\nonumber\\
&  \left[  q_{12},q_{23}\right]  =\left[  c_{1}a_{1}^{\dagger}a_{2}%
-c_{2}a_{-2}^{\dagger}a_{-1},c_{2}a_{2}^{\dagger}a_{3}-c_{3}a_{-3}^{\dagger
}a_{-2}\right]  =\nonumber\\
&  c_{2}\left(  c_{1}a_{1}^{\dagger}a_{3}-c_{3}a_{-3}^{\dagger}a_{-1}\right)
=c_{2}q_{13}\nonumber\\
&  \left[  q_{12},q_{2-3}\right]  =\left[  c_{1}a_{1}^{\dagger}a_{2}%
-c_{2}a_{-2}^{\dagger}a_{-1},c_{2}a_{2}^{\dagger}a_{-3}-c_{3}a_{3}^{\dagger
}a_{-2}\right]  =\nonumber\\
&  -c_{2}\left(  c_{1}a_{1}^{\dagger}a_{-3}+c_{3}a_{3}^{\dagger}a_{-1}\right)
=-c_{2}q_{1-3}\nonumber\\
&  \left[  q_{12},q_{-23}\right]  =\left[  c_{1}a_{1}^{\dagger}a_{2}%
-c_{2}a_{-2}^{\dagger}a_{-1},c_{2}a_{-2}^{\dagger}a_{3}-c_{3}a_{-3}^{\dagger
}a_{2}\right]  =0\nonumber\\
&  \left[  q_{12},q_{-2-3}\right]  =\left[  c_{1}a_{1}^{\dagger}a_{2}%
-c_{2}a_{-2}^{\dagger}a_{-1},c_{2}a_{-2}^{\dagger}a_{-3}-c_{3}a_{3}^{\dagger
}a_{2}\right]  =0
\end{align}


which covers and corresponds to $-3\leq j\leq k\leq3.$

\subsection{Lie Algebra Generators}

A reducible represenation $R\left(  gl\left(  2s,\mathbb{C}\right)  \right)  $
of the Lie algebra $gl\left(  2s,\mathbb{C}\right)  $ on $\mathcal{H}_{F}$ is
generated by%

\begin{align}
q_{jk}^{\dagger}  &  =c_{j}a_{k}^{\dagger}a_{j}-\sigma\left(  j\right)
\sigma\left(  k\right)  c_{k}a_{-j}^{\dagger}a_{-k};1\leq j<k\leq s;\left\vert
j\right\vert \neq\left\vert k\right\vert \nonumber\\
q_{jk}  &  =c_{j}a_{j}^{\dagger}a_{k}-\sigma\left(  j\right)  \sigma\left(
k\right)  c_{k}a_{-k}^{\dagger}a_{-j};1\leq j<k\leq s;\left\vert j\right\vert
\neq\left\vert k\right\vert \nonumber\\
a_{j}^{\dagger}a_{k};\left\vert j\right\vert  &  =\left\vert k\right\vert
;1\leq j\leq s.
\end{align}
In line with Gilmore (where $q\equiv Z$ ) one defines
\begin{align}
q_{j-j}  &  =a_{j}^{\dagger}a_{-j}\nonumber\\
q_{-jj}  &  =a_{-j}^{\dagger}a_{j}=q_{j-j}^{\dagger}.
\end{align}
The number of $q$ generators these generators is
\begin{equation}
2s(2s-1)-2=4s^{2}-4s
\end{equation}
and the number of subblock generators is%

\begin{equation}
4s
\end{equation}
giving a total number
\begin{equation}
4s^{2},
\end{equation}
which equals the number of generators $\left\{  \left.  a_{j}^{\dagger}%
a_{k}\right\vert -s\leq j,k\leq s\right\}  .$ The representation is the sum of
the irreps
\begin{equation}
R\left(  gl\left(  2s,\mathbb{C}\right)  \right)  =\bigoplus\limits_{0\leq
N\leq2s}R_{N}\left(  gl\left(  2s,\mathbb{C}\right)  \right)  ,
\end{equation}
where $\left\{  \left.  R_{N}\left(  gl\left(  2s,\mathbb{C}\right)  \right)
\right\vert 0\leq N\leq r\right\}  $ are the irreducible representions carried
by $\mathcal{H}^{N}.$

The $\left\{  q_{jk},q_{jk}^{\dagger}\left\vert 1\leq j<k\leq s;\left\vert
j\right\vert \neq\left\vert k\right\vert \right.  \right\}  $ generators can
be expressed in terms of $\left\{  \left.  a_{j}^{\dagger}a_{k}\right\vert
-s\leq j<k\leq s;\left\vert j\right\vert \neq\left\vert k\right\vert \right\}
,$ when $c_{j}\neq c_{k},$ as
\begin{equation}
\left(
\begin{array}
[c]{c}%
q_{-j-k}^{\dagger}\\
q_{j-k}^{\dagger}\\
q_{-jk}^{\dagger}\\
q_{jk}^{\dagger}\\
q_{-j-k}\\
q_{j-k}\\
q_{-jk}\\
q_{jk}%
\end{array}
\right)  =\left(
\begin{array}
[c]{cccccccc}%
c_{k} & 0 & \cdots & \cdots & \cdots & \cdots & 0 & -c_{j}\\
0 & c_{k} & 0 & \cdots & \cdots & 0 & c_{j} & 0\\
\vdots & 0 & c_{k} & 0 & 0 & c_{j} & 0 & \vdots\\
\vdots & \vdots & 0 & c_{k} & -c_{j} & 0 & \vdots & \vdots\\
\vdots & \vdots & 0 & -c_{j} & c_{k} & 0 & \vdots & \vdots\\
\vdots & 0 & c_{j} & 0 & 0 & c_{k} & 0 & \vdots\\
0 & c_{j} & 0 & \cdots & \cdots & 0 & c_{k} & 0\\
-c_{j} & 0 & \cdots & \cdots & \cdots & \cdots & 0 & c_{k}%
\end{array}
\right)  \left(
\begin{array}
[c]{c}%
a_{-j}^{\dagger}a_{-k}\\
a_{-j}^{\dagger}a_{k}\\
a_{j}^{\dagger}a_{-k}\\
a_{j}^{\dagger}a_{k}\\
a_{-k}^{\dagger}a_{-j}\\
a_{k}^{\dagger}a_{-j}\\
a_{-k}^{\dagger}a_{j}\\
a_{k}^{\dagger}a_{j}%
\end{array}
\right)  .
\end{equation}%
\begin{equation}
\left(
\begin{array}
[c]{c}%
u_{jk}\\
u_{j+sk}\\
u_{jk+s}\\
u_{j+sk+s}\\
v_{jk}\\
v_{j+sk}\\
v_{jk+s}\\
v_{j+sk+s}%
\end{array}
\right)  =\left(
\begin{array}
[c]{cccccccc}%
c_{k} & 0 & \cdots & \cdots & \cdots & \cdots & 0 & c_{j}\\
0 & c_{k} & 0 & \cdots & \cdots & 0 & -c_{j} & 0\\
\vdots & 0 & c_{k} & 0 & 0 & -c_{j} & 0 & \vdots\\
\vdots & \vdots & 0 & c_{k} & c_{j} & 0 & \vdots & \vdots\\
\vdots & \vdots & 0 & -c_{j} & c_{k} & 0 & \vdots & \vdots\\
\vdots & 0 & c_{j} & 0 & 0 & c_{k} & 0 & \vdots\\
0 & c_{j} & 0 & \cdots & \cdots & 0 & c_{k} & 0\\
-c_{j} & 0 & \cdots & \cdots & \cdots & \cdots & 0 & c_{k}%
\end{array}
\right)  \left(
\begin{array}
[c]{c}%
a_{j}^{\dagger}a_{k}\\
a_{j}^{\dagger}a_{k+s}\\
a_{j+s}^{\dagger}a_{k}\\
a_{j+s}^{\dagger}a_{k+s}\\
a_{k}^{\dagger}a_{j}\\
a_{k}^{\dagger}a_{j+s}\\
a_{k+s}^{\dagger}a_{j}\\
a_{k+s}^{\dagger}a_{j+s}%
\end{array}
\right)  ,
\end{equation}%
\begin{align}
\left[  \varkappa_{jk},\varkappa_{lm}\right]   &  =\left[  \sum
_{\substack{-s\leq j,k\leq s,\left\vert j\right\vert \neq\left\vert
k\right\vert \\1\leq\alpha,\beta\leq8}}c_{jk\alpha\beta}a_{\alpha}^{\dagger
}a_{\beta},\sum_{\substack{-s\leq l,m\leq s,\left\vert l\right\vert
\neq\left\vert m\right\vert \\1\leq\gamma,\delta\leq8}}c_{lm\gamma\delta
}a_{\gamma}^{\dagger}a_{\delta}\right] \\
&  =\sum_{\substack{-s\leq j,k\leq s,\left\vert j\right\vert \neq\left\vert
k\right\vert \\1\leq\alpha,\beta\leq8\\-s\leq l,m\leq s,\left\vert
l\right\vert \neq\left\vert m\right\vert \\1\leq\gamma,\delta\leq
8}}c_{jk\alpha\beta}c_{lm\gamma\delta}\left[  a_{\alpha}^{\dagger}a_{\beta
},a_{\gamma}^{\dagger}a_{\delta}\right]
\end{align}


\section{Commutation Relationships}

The commutation relations for $sp\left(  2s,\mathbb{C}\right)  $ are
\begin{align}
\left[  q_{ij},q_{kl}\right]   &  =\sigma\left(  jk\right)  \left\{
q_{il}\delta_{kj}+q_{-j-k}\delta_{il}+q_{i-k}\delta_{-jl}+q_{-jl}\delta
_{k-i}\right\}  ;\\
-s  &  \leq i<j\leq s;-s\leq k<l\leq s
\end{align}
and the symmetry relationships
\begin{equation}
q_{ij}=-\sigma\left(  ij\right)  q_{-j-i}%
\end{equation}
are satisfied.%

\begin{equation}
\left[  a_{r}^{\dagger}a_{s},a_{p}^{\dagger}a_{q}\right]  =a_{r}^{\dagger
}a_{q}\delta_{sp}-a_{p}^{\dagger}a_{s}\delta_{rq}%
\end{equation}%
\begin{align}
\left[  a_{r}^{\dagger}a_{s},a_{p}^{\dagger}a_{p}\right]   &  =a_{r}^{\dagger
}a_{p}\delta_{sp}-a_{p}^{\dagger}a_{s}\delta_{rp}\\
&  =\left\{
\begin{array}
[c]{c}%
a_{r}^{\dagger}a_{p},s=p,r\neq s\\
-a_{p}^{\dagger}a_{s},r=p,r\neq s
\end{array}
\right.
\end{align}%
\begin{align}
&  \left[  \left(  c_{j}a_{k}^{\dagger}a_{j}-\sigma\left(  j\right)
\sigma\left(  k\right)  c_{k}a_{-j}^{\dagger}a_{-k}\right)  ,a_{l}^{\dagger
}a_{m}\right]  =\\
&  c_{j}\left(  a_{k}^{\dagger}a_{m}\delta_{jl}-a_{l}^{\dagger}a_{j}%
\delta_{km}\right)  -c_{k}\sigma\left(  j\right)  \sigma\left(  k\right)
\left(  a_{-j}^{\dagger}a_{m}\delta_{-kl}-a_{l}^{\dagger}a_{-k}\delta
_{-jm}\right)
\end{align}%
\begin{equation}
m=l=j:\left[  ,\right]  =c_{j}a_{k}^{\dagger}a_{m}-c_{k}\sigma\left(
j\right)  \sigma\left(  k\right)
\end{equation}


If one shifts the indexing by
\begin{align}
\upsilon\left(  j\right)   &  =2j,1\leq j\leq s\\
\upsilon\left(  j\right)   &  =-\left(  2j+1\right)  ,-s\leq j\leq-1
\end{align}
then%
\begin{equation}
\left[  q_{ij},q_{jk}\right]  =q_{ik};1\leq i<j<k\leq2s
\end{equation}
and
\begin{equation}
\left[  q_{jk}^{\dagger},q_{ij}^{\dagger}\right]  =-\left[  q_{ij}^{\dagger
},q_{jk}^{\dagger}\right]  =q_{ik}^{\dagger};1\leq i<j<k\leq2s.
\end{equation}
The general linear group $GL\left(  \mathcal{H}^{1}\right)  $ can be factored
into the product%
\begin{equation}
GL\left(  \mathcal{H}^{1}\right)  =GL\left(  1\right)  \times SL\left(
\mathcal{H}^{1}\right)  ,
\end{equation}
such that
\begin{equation}
\mathrm{g}=\left(  \mathrm{\det g}\right)  \mathrm{h}=e^{Tr\left\{
\mathrm{g}\right\}  }\mathrm{h,h\in}SL\left(  \mathcal{H}\right)  .
\end{equation}
where
\begin{equation}
\eta_{j}=Tr\left\{  \mathbf{\lambda}_{j}\right\}
\end{equation}


\section{}
\end{document}